%!TEX program = xelatex
%!TEX root = ../all.tex

\chapter{Support Vector Machines}

\section*{Idea}

Support Vector Machines (SVMs) address the problem of binary classification by searching for a hyperplane that separates two classes in a given feature space. The central idea, however, is not merely to find \emph{any} separating hyperplane, but to identify one that is optimal according to a precise and geometrically meaningful criterion: the maximization of the margin. Intuitively, among all possible decision boundaries that correctly classify the training data, we prefer the one that leaves the largest possible ``safety buffer'' between the two classes.

This perspective immediately highlights two fundamental difficulties that arise in practice. First, real data are often not perfectly linearly separable, so insisting on flawless separation may be unrealistic or even harmful. Second, even when linear separation fails in the original feature space, it may become possible after a suitable transformation of the data. These considerations motivate two complementary extensions of the basic idea. On the one hand, we can relax the notion of separation by allowing some violations of the margin constraints, leading to the concept of \emph{soft margins}. On the other hand, we can enrich the feature space, typically through nonlinear mappings implemented via kernel functions, so that a linear separator in the new space corresponds to a nonlinear decision boundary in the original one.

Before formalizing these ideas, it is useful to recall how linear classifiers produce predictions and how these predictions are connected to geometry.

\subsection*{Classification Scores and Decision Boundaries}

Most classification algorithms do not directly output a class label, but rather a numerical score that quantifies how strongly a data point is associated with one class or the other. The final label is then obtained by comparing this score with a fixed threshold. In the case of linear classifiers, this score is an affine function of the input features:
\begin{myequation*}
\text{score}(\x)=\w^T\x+b .
\end{myequation*}

The vector $\w$ determines the orientation of the decision boundary, while the scalar $b$ controls its position. Different learning algorithms interpret this score in different ways. For instance, logistic regression maps the score through a sigmoid function to produce an estimate of the conditional class probability,
\begin{myequation*}
p(y=1\mid\x)=\sigma(\w^T\x+b)\in[0,1],
\end{myequation*}
and assigns the label $1$ whenever this probability exceeds $0.5$. Other methods, such as the perceptron or the hard-margin SVM, use only the sign of the score and do not attempt to estimate probabilities.

Beyond its algorithmic role, the score has a clear geometric interpretation. The set of points satisfying $\w^T\x+b=0$ forms a hyperplane, known as the \emph{decision boundary}. The value of $\w^T\x+b$ is proportional to the signed distance of $\x$ from this hyperplane: its sign indicates on which side of the boundary the point lies, while its magnitude reflects how far the point is from the boundary. Points that are far away are classified with high confidence, whereas points close to the hyperplane are inherently ambiguous.

\begin{center}
\begin{tikzpicture}[
    thick,
    >=stealth',
    dot/.style = {
      draw,
      fill = white,
      circle,
      inner sep = 0pt,
      minimum size = 4pt
    }
  ]
  \coordinate (O) at (0,0);
  \draw[->] (-0.3,0) -- (8,0) coordinate[label = {below:$w_1$}] (xmax);
  \draw[->] (0,-0.3) -- (0,5) coordinate[label = {right:$w_2$}] (ymax);
  \path[name path=x] (0.3,0.5) -- (6.7,4.7);
  \path[name path=y] plot[smooth] coordinates {(-0.3,2) (2,1.5) (4,2.8) (6,5)};
  \scope[name intersections = {of = x and y, name = i}]
    \draw[red] (-1,4.66) -- (7.5,-1) node[pos=0.7,  right] {\color{black}\footnotesize$\w^T\x-b=0$};
     \fill (5,3) circle (1.5pt);
    \node[right] at (5,3) {$\mathbf{\hat x}$};
     \draw[blue, ->] (0,0) --  (5,3);
      \fill (2.92,4.38) circle (1.5pt);
      \draw[gray,dashed](5,3)--(2.92,4.38);
    \draw[red, ->] (0,0) -- node[right] {$\w$} (1,1.5);
    \draw[red, dashed] (0,0) -- (3.5,5.25);
    
     \fill (-.15,.1) circle (.3pt);
     \fill (2.77,4.48) circle (.3pt);
     \draw[gray, line width=.5pt] (-.225,.15) -- (0,0);
     \draw[gray, line width=.5pt] (-.15,.1) -- (2.77,4.48);
     \draw[gray, line width=.5pt] (2.92,4.38) -- (2.695,4.53);
     \node[gray] at (.5,2) {$\w^T\mathbf{\hat x}$};
      \draw[gray, line width=.5pt] (-1.5,.99) -- (0,0);
       \fill (.5,3.66) circle (1pt);
       \fill (-1.34,.89) circle (1pt);
        \draw[gray, line width=.5pt] (.5,3.66) -- (-1.34,.89);
        \node[gray] at (-.7,2.5) {$-b$};
         \draw[gray, line width=1pt] (5,3) -- (3.92,1.38);
          \node[gray] at (5.2,2) {$\w^T\mathbf{\hat x}+b$};
  \endscope
\end{tikzpicture}
\end{center}

This geometric viewpoint is illustrated further by comparing points at different distances from the decision boundary. In the figure below, point $A$ lies far from the separating hyperplane and is therefore classified with high confidence, whereas point $C$, being very close to the boundary, corresponds to a much more uncertain decision. Such observations naturally suggest that, all else being equal, a classifier that places training points far from the decision boundary should generalize better to unseen data.

\centerline{
\includegraphics[width=.4\linewidth]{largemargin.pdf}
}

\subsubsection*{Margin: Formal Definition}

To make these ideas precise, consider a binary classification problem with labels $y\in\{-1,1\}$, where $y=1$ denotes class $\mathcal{C}_1$ and $y=-1$ denotes class $\mathcal{C}_0$. We restrict our attention to linear classifiers of the form
\begin{myequation*}
h(\x)=\text{sign}(\w^T\x+b),
\end{myequation*}
which return only a class label and no probabilistic information.

For a training example $(\x_i,t_i)$, with true label $t_i\in\{-1,1\}$, we define the \alert{functional margin} as
\begin{myequation*}
\overline{\gamma}_i = t_i(\w^T\x_i+b).
\end{myequation*}
This quantity is positive if and only if the example is correctly classified, and its magnitude increases as the point moves away from the decision boundary in the correct direction. From this point of view, the functional margin can be interpreted as a signed confidence measure associated with the prediction on $\x_i$.

Extending this notion to the entire training set $\tset=\{(\x_1,t_1),\ldots,(\x_n,t_n)\}$, we define the functional margin of the classifier as the minimum functional margin over all examples,
\begin{myequation*}
\overline{\gamma}=\min_{i}\overline{\gamma}_i.
\end{myequation*}
Maximizing this quantity would seem to encourage confident predictions on all training points. However, the functional margin suffers from a fundamental limitation: it is not invariant under rescaling of the parameters. Multiplying $(\w,b)$ by any positive constant $\alpha$ leaves the classifier unchanged, but multiplies $\overline{\gamma}_i$ by $\alpha$. As a result, the functional margin alone cannot serve as an intrinsic measure of geometric separation.

This issue is resolved by introducing the \alert{geometric margin}, which directly measures distances in the input space. For a single example, the geometric margin is defined as
\begin{myequation*}
\gamma_i = t_i\left(\frac{\w^T}{\|\w\|}\x_i + \frac{b}{\|\w\|}\right)
= \frac{\overline{\gamma}_i}{\|\w\|},
\end{myequation*}
where $\|\w\|$ denotes the Euclidean norm of $\w$. By construction, $\gamma_i$ is exactly the signed Euclidean distance from $\x_i$ to the separating hyperplane. Importantly, this quantity is invariant under rescaling of the parameters, and thus captures a true geometric property of the classifier.

\begin{center}
\includegraphics[width=.4\linewidth]{margin.pdf}
\end{center}

The geometric margin of the entire training set is then defined as
\begin{myequation*}
\gamma = \min_{i} \gamma_i.
\end{myequation*}
Geometrically, $\gamma$ corresponds to half the width of the largest slab, centered on the decision hyperplane, that contains no training points. The two hyperplanes delimiting this slab are parallel to the decision boundary and are located at distance $\gamma$ from it. These are known as the \alert{maximum margin hyperplanes}. Training points that lie exactly on these boundaries play a crucial role and are referred to as \alert{support vectors}.

\begin{center}
 \begin{tikzpicture}{0cm}{-3cm}{5cm}{5cm}
\draw[->] (-.5,0) -- (10,0) node[right] {};
\draw[->] (0,-.5) -- (0,5) node[above] {};
\draw[-] (-1,4.5) -- (10,-1) node[above] {};
\draw[dashed] (-1,5) -- (11,-1) node[above] {};
\draw[dashed] (-1,4) -- (9,-1) node[above] {};
\draw[<->] (1,3) -- (1.4,3.8) node[left] {};
\node at (1.6,3.5) {\small $2\gamma$};
\node at (5,3) {{\tiny x}};
\node at (4,4.5) {{\tiny x}};
\node at (6,1.7) {{\tiny x}};
\node at (6.5,5.3) {{\tiny x}};
\node at (7.2,3.9) {{\tiny x}};
\node at (7.4,1) {{\tiny x}};
\node at (1.2,4) {{\tiny x}};
\node at (2,3.9) {{\tiny x}};
\node at (3,3.8) {{\tiny x}};
\node at (3.8,2.7) {{\tiny x}};
\node at (4.2,2.4) {{\tiny x}};
\filldraw [black] (1,0.5) circle (1pt)
(0.8,1.3) circle (1pt)
(2,1.5) circle (1pt)
(3.8,0.5) circle (1pt)
(3.4,0.7) circle (1pt)
(3.5,0.3) circle (1pt)
(0.2,3) circle (1pt)
(1.7,2.1) circle (1pt)
(4.7,0.3) circle (1pt)
(2.7,0.8) circle (1pt)
(6.5,0.25) circle (1pt);
\pgfusepath{fill}
\end{tikzpicture}
\end{center}

\subsubsection*{Interpreting Margin Constraints}

The value of $t(\w^T\x+b)$ provides a detailed description of the position of a point relative to the decision boundary and the margin strip. Values larger than $1$ correspond to correctly classified points that lie outside the margin, indicating high confidence predictions. When the value is exactly $1$, the point lies on the margin boundary and is a support vector. Values between $0$ and $1$ indicate correctly classified points that lie inside the margin, hence with lower confidence. Negative values correspond to misclassified points, with the magnitude again reflecting how far the point lies on the wrong side of the boundary. This continuum of situations is essential for understanding both hard- and soft-margin formulations.

\section*{Optimal Margin Classifiers}

Armed with these definitions, we can now formalize the problem of finding an optimal separating hyperplane. When the training data are linearly separable, the goal is to find the classifier that maximizes the geometric margin. From a statistical perspective, a large margin acts as a form of regularization, reducing the sensitivity of the classifier to noise and small perturbations in the data, and thereby improving generalization performance.

\subsection*{The Optimization Problem}

For a linearly separable training set, the problem can be stated as
\begin{myalign}
&\max_{\w,b} \ \gamma \\
&\text{subject to} \quad \gamma_i = \frac{t_i}{\|\w\|}(\w^T\x_i+b) \ge \gamma,
\quad i=1,\ldots,n .
\end{myalign}
Multiplying both sides of the constraints by $\|\w\|$ yields an equivalent formulation,
\begin{myalign}
&\max_{\w,b} \ \gamma \\
&\text{subject to} \quad t_i(\w^T\x_i+b) \ge \gamma \|\w\|,
\quad i=1,\ldots,n .
\end{myalign}

The key observation is that the decision boundary is invariant under positive rescaling of $(\w,b)$. We can therefore fix this degree of freedom by imposing the normalization condition
\begin{myequation*}
\min_{i} t_i(\w^T\x_i+b)=1.
\end{myequation*}
Under this constraint, the margin becomes $\gamma = 1/\|\w\|$, and maximizing the margin is equivalent to minimizing the norm of $\w$. This leads to the canonical hard-margin SVM optimization problem,
\begin{myalign}
&\min_{\w,b} \ \frac{1}{2}\|\w\|^2 \\
&\text{subject to} \quad t_i(\w^T\x_i+b) \ge 1,
\quad i=1,\ldots,n .
\end{myalign}

%\subsubsection*{Properties of the Optimization Problem}

The resulting optimization problem is a convex quadratic program. The objective function is strictly convex in $\w$, and the constraints define a convex feasible region given by the intersection of half-spaces. As a consequence, the problem admits a unique global optimum, which can be found efficiently using well-established optimization techniques.

At the optimum, not all constraints are equally important. Points for which
\begin{myequation*}
t_i(\w^T\x_i+b)=1
\end{myequation*}
are said to be \alert{active} and correspond precisely to the support vectors. These points alone determine the position of the optimal hyperplane. All other training examples, which satisfy the constraints strictly, could be removed without affecting the solution. This sparsity property is one of the most appealing features of SVMs.

\subsection*{Duality in Convex Optimization}

A fundamental property of convex optimization problems is the existence of an equivalent \textit{dual} formulation. The dual problem provides an alternative perspective on the same optimization task and, under certain conditions, yields computational advantages. This section develops the concept of duality rigorously.

\subsubsection*{Problem Setup}

Consider the constrained optimization problem:
\begin{myalign}
&\min_{\x \in \Omega} f(\x)
\end{myalign}
where the feasible region $\Omega$ is defined by inequality constraints:
\begin{myalign}
&g_i(\x) \leq 0 \quad \text{for } i=1,\ldots,k
\end{myalign}

Throughout this section, we assume that $f(\x)$ and all $g_i(\x)$ are convex functions, and that $\Omega$ is a convex set. We refer to this as the \alert{primal problem}.

\subsubsection*{The Lagrangian}

The notion of a \alert{Lagrangian} arises naturally when we try to formalize the idea of optimization in the presence of constraints. In many practical problems, we are not free to choose the variables $\x$ arbitrarily: instead, admissible solutions must satisfy a collection of constraints that encode physical limitations, resource bounds, consistency requirements, or modeling assumptions. The central conceptual difficulty is that constrained optimization couples two competing objectives: on the one hand, we wish to minimize (or maximize) a given cost function; on the other hand, we must enforce that the solution remains within the feasible region defined by the constraints.

A useful way to think about this situation is to imagine that constraint violations come with a \emph{cost}. Rather than treating constraints as rigid, binary conditions (either satisfied or violated), we temporarily allow violations but penalize them in a controlled way. The Lagrangian provides exactly this mechanism: it transforms a constrained optimization problem into an unconstrained one by augmenting the original objective with penalty terms that measure how much each constraint is violated.

More precisely, consider an optimization problem in which we aim to minimize a function $f(\x)$ subject to inequality constraints
\begin{myequation*}
g_i(\x) \le 0, \qquad i = 1, \ldots, k.
\end{myequation*}
The Lagrangian associated with this problem is defined as
\begin{myequation*}
L(\x, \Vlambda) = f(\x) + \sum_{i=1}^k \lambda_i g_i(\x),
\end{myequation*}
where the vector $\Vlambda = (\lambda_1, \ldots, \lambda_k)$ collects the \alert{Lagrangian multipliers}, also known as \alert{dual variables}. Each multiplier $\lambda_i$, which is bound to be non negative, is associated with a specific constraint $g_i(\x) \le 0$ and controls how strongly violations of that constraint are penalized.

This construction admits a clear and intuitive interpretation. If, for a given $\x$, a constraint $g_i(\x)$ is satisfied, then $g_i(\x) \le 0$ and the corresponding term $\lambda_i g_i(\x)$ does not increase the value of the Lagrangian. If instead the constraint is violated, $g_i(\x) > 0$, the penalty term becomes positive and increases $L(\x,\Vlambda)$ in proportion both to the amount of violation and to the multiplier $\lambda_i$. In this sense, the multipliers play the role of adjustable weights that determine how costly it is to violate each constraint.

From a geometric perspective, the Lagrangian blends information about the objective function and the constraints into a single scalar quantity. Minimizing $L(\x,\Vlambda)$ with respect to $\x$ for fixed $\Vlambda$ corresponds to searching for solutions that strike a balance between reducing the objective $f(\x)$ and avoiding constraint violations. At the same time, varying $\Vlambda$ changes this balance, effectively reshaping the optimization landscape so as to enforce feasibility.

Each multiplier $\lambda_i$ can also be viewed as a \emph{shadow price}: it quantifies how much the optimal value of the objective would improve if the corresponding constraint were relaxed slightly. In this sense, the dual variables encode the sensitivity of the optimal solution to the constraints and provide valuable insight beyond the primal solution $\x$ itself.
\medskip

\subsubsection*{Weak Duality}

Starting from the Lagrangian, duality can be introduced in a particularly intuitive way by thinking in terms of \emph{two players} and \emph{two complementary viewpoints} on the same constrained optimization problem.

\medskip

%Suppose we wish to minimize an objective $f(\x)$ over a feasible set defined by constraints. In the \alert{primal} viewpoint, $\x$ is the decision variable and feasibility is essential: we want the best value of $f(\x)$ among admissible solutions. The Lagrangian allows us to encode constraints by augmenting the objective with weighted violations. This immediately suggests an adversarial interpretation: if $\x$ violates a constraint, a multiplier $\lambda_i$ can be increased to penalize that violation, making the choice of $\x$ unattractive. In other words, the multipliers can be used to \emph{enforce feasibility} indirectly.
%
%This leads naturally to the following minimax expression. For a fixed $\x$, consider maximizing the Lagrangian over multipliers $\Vlambda$ with $\lambda_i \ge 0$. If $\x$ is feasible, then each $g_i(\x)\le 0$ and increasing $\lambda_i$ cannot increase $L(\x,\Vlambda)$, so the maximization does not ``punish'' $\x$. If $\x$ is infeasible and some constraint has $g_i(\x)>0$, then the maximization drives $\lambda_i$ upward and makes $L(\x,\Vlambda)$ arbitrarily large. Hence the maximization over $\Vlambda$ acts as a mechanism that assigns an infinite cost to infeasible points and leaves feasible points essentially unchanged. Minimizing over $\x$ after this ``penalty maximization'' reproduces the original constrained problem. 

The key observation is that maximizing the Lagrangian over $\Vlambda \ge 0$ for a fixed
$\x$ recovers exactly the primal cost:
\begin{myequation*}
\max_{\Vlambda \ge 0} L(\x, \Vlambda) \;=\;
\begin{cases}
f(\x) & \text{if } \x \text{ is feasible, i.e., } g_i(\x) \le 0 \;\forall\, i, \\
+\infty & \text{if } \x \text{ is infeasible, i.e., } g_i(\x) > 0 \text{ for some } i.
\end{cases}
\end{myequation*}
To see why, note that if $\x$ is feasible then all $g_i(\x) \le 0$, so each term
$\lambda_i g_i(\x) \le 0$, and the maximum over $\Vlambda \ge 0$ is achieved by setting all
$\lambda_i = 0$, giving $L(\x, \Vlambda) = f(\x)$. If instead some constraint is violated,
say $g_j(\x) > 0$, then sending $\lambda_j \to +\infty$ drives $L(\x, \Vlambda) \to
+\infty$. The maximization over $\Vlambda$ therefore acts as an automatic feasibility
enforcer: it assigns an infinite cost to any infeasible point and leaves the cost of feasible
points unchanged.

The \alert{primal function}
\begin{myequation*}
p(\x) \;=\; \max_{\Vlambda \ge 0}\, L(\x, \Vlambda) \;=\; \max_{\Vlambda \ge 0}\left[f(\x) +
\sum_{i=1}^k \lambda_i g_i(\x)\right].
\end{myequation*}
associates to each $\x$ a real value $p(\x)$, which is the largest value the Lagrangian can take, maximized over all multipliers, bound to be non negative. As observed, $p(\x)$ is infinite if and only if $\x$ is an unfeasible point.

The primal problem asks for the minimization of $p(\x)$ and is thus exactly equivalent to a minimax problem
\begin{myequation*}
p^* \;=\; \min_{\x}\, p(\x) \;=\; \min_{\x}\, \max_{\Vlambda \ge 0}\, L(\x, \Vlambda),
\end{myequation*}
where no explicit feasibility constraint on $\x$ appears: feasibility is enforced implicitly
through the maximization of the multipliers.



\subsubsection*{The dual problem.}
Duality arises by swapping the order of minimization and maximization. Define the
\alert{dual function} as
\begin{myequation*}
q(\Vlambda) \;=\; \min_{\x}\, L(\x, \Vlambda) \;=\; \min_{\x}\left[f(\x) +
\sum_{i=1}^k \lambda_i g_i(\x)\right].
\end{myequation*}
This clearly associates a real value $q(\Vlambda)$ to each $\Vlambda \ge 0$, which is the smallest value the Lagrangian can
take, minimized freely over $\x$ without any feasibility constraint. The \alert{dual
problem} is then the maximization of this function and is equivalent to a maxmin problem
\begin{myequation*}
d^* \;=\; \max_{\Vlambda \ge 0}\, q(\Vlambda) \;=\; \max_{\Vlambda \ge 0}\, \min_{\x}\, L(\x, \Vlambda),
\end{myequation*}
Geometrically, $q(\Vlambda)$ is a lower bound on $p^*$ for every $\Vlambda \ge 0$. 
To see why this holds, consider any feasible point
$\x$: since in this case $\lambda_i \ge 0$
and $g_i(\x) \le 0$, every penalty term satisfies $\lambda_i g_i(\x) \le 0$, and therefore
\begin{myequation*}
L(\x, \Vlambda) \;=\; f(\x) + \sum_{i=1}^k \lambda_i g_i(\x) \;\le\; f(\x).
\end{myequation*}
In other words, for any feasible $\x$ and any $\Vlambda \ge 0$, the Lagrangian is never
larger than the original objective. Taking the minimum over $\x$ on both sides, and note
that the left-hand side minimizes over \emph{all} $\x$ (not just feasible ones), while the
right-hand side minimizes only over feasible $\x$. This results into
\begin{myequation*}
q(\Vlambda) \;=\; \min_{\x}\, L(\x, \Vlambda) \;\le\; \min_{\x \text{ feasible}}\,
L(\x, \Vlambda) \;\le\; \min_{\x \text{ feasible}}\, f(\x) \;=\; p^*.
\end{myequation*}
The first inequality holds because minimizing over a larger set can only decrease the
minimum. The second inequality follows from the bound $L(\x, \Vlambda) \le f(\x)$
established above, which holds for every feasible $\x$.

Since this chain of inequalities holds for every $\Vlambda \ge 0$, it holds in particular
for the dual maximizer $\Vlambda^*$, giving $q(\Vlambda^*) \le p^*$. This gives the fundamental \alert{weak duality} inequality
\begin{myequation*}
d^* \;\le\; p^*,
\end{myequation*}
which holds always, regardless of convexity or any other structural assumption on $f$ and $g_i$.
The gap $p^* -
d^* \ge 0$ is called the \alert{duality gap}, and it measures how much information is lost
by swapping the order of minimization and maximization.


%\subsubsection*{Strong duality.}
%Under suitable conditions --- most notably when $f$ and the $g_i$ are convex and a strictly
%feasible point exists (Slater's condition) --- the duality gap closes and
%\begin{myequation*}
%d^* \;=\; p^*.
%\end{myequation*}
%This is \alert{strong duality}. When it holds, solving the dual problem is equivalent to
%solving the primal: one can recover the primal optimal $\x^*$ from the dual optimal
%$\Vlambda^*$. The hard-margin SVM objective is a convex quadratic program with linear
%constraints, so Slater's condition is satisfied whenever the data are linearly separable, and
%strong duality holds. This justifies the derivation of the dual SVM formulation, to which we
%turn next.
%
%
%This is the primal formulation written in terms of the Lagrangian:
%\begin{myalign}
%p^* 
%= \min_{\x \in \Omega} f(\x)
%= \min_{\x \in \Omega} \max_{\Vlambda:\,\lambda_i \ge 0} L(\x,\Vlambda).
%\end{myalign}
%Here, $\Omega$ denotes the domain of $\x$ (for example $\mathbb{R}^d$ or another set encoding basic restrictions), while the constraints are enforced through the maximization over $\Vlambda$.
%
%\medskip
%
%The \alert{dual} viewpoint reverses the order of these operations. Instead of choosing $\x$ first, we fix multipliers $\Vlambda$ and ask: \emph{what is the smallest value of the Lagrangian we can obtain over $\x$?} This gives rise to the \alert{dual function}
%\begin{myequation*}
%q(\Vlambda) = \min_{\x\in\Omega} L(\x,\Vlambda),
%\end{myequation*}
%which assigns to each set of multipliers the best (i.e., lowest) value of the penalized objective that can be achieved. Importantly, for any $\Vlambda \succeq 0$ this quantity provides a \emph{lower bound} on the primal optimum $P^*$ (for minimization problems). The intuition is the following: when $\x$ is feasible, $g_i(\x)\le 0$ and thus $\sum_i \lambda_i g_i(\x)\le 0$, meaning that
%\begin{myequation*}
%L(\x,\Vlambda) = f(\x) + \sum_i \lambda_i g_i(\x) \le f(\x).
%\end{myequation*}
%Taking the minimum over $\x$ on the left-hand side cannot exceed the minimum of $f(\x)$ over feasible points, so $q(\Vlambda)\le P^*$ for all admissible $\Vlambda$. Therefore, by searching over multipliers we can try to make this bound as tight as possible.
%
%This idea defines the dual optimization problem:
%\begin{myalign}
%D^* 
%= \max_{\Vlambda:\,\lambda_i \ge 0} \min_{\x \in \Omega} L(\x, \Vlambda)
%= \max_{\Vlambda:\,\lambda_i \ge 0} q(\Vlambda).
%\end{myalign}
%In words, the dual problem chooses the multipliers so that the \emph{best lower bound} produced by the Lagrangian is as large as possible.
%
%
%\subsubsection*{Relationship between primal and dual solutions.}
%
%The connection between primal and dual solutions can be understood most clearly by interpreting the Lagrangian as a device that simultaneously evaluates the objective function and measures constraint violations. From this perspective, primal and dual variables play complementary roles: primal variables propose candidate solutions, while dual variables assess and penalize their feasibility.
%
%\medskip
%
%A central consequence of this interaction is the following fundamental principle:
%
%\medskip
%\noindent
%\emph{Any feasible dual solution provides a lower bound on the optimal value of the primal problem.}
%
%\medskip
%
%This statement is the essence of \alert{weak duality}. To understand why it holds, it is instructive to analyze how the Lagrangian behaves depending on whether a candidate primal point $\x$ satisfies the constraints or not.
%
%\medskip
%
%Consider first a point $\x$ that is \textbf{feasible}, meaning that all inequality constraints are satisfied:
%\begin{myequation*}
%g_i(\x) \leq 0 \qquad \forall i.
%\end{myequation*}
%For any choice of multipliers $\lambda_i \geq 0$, each product $\lambda_i g_i(\x)$ is non-positive. As a consequence, the Lagrangian satisfies
%\begin{myequation*}
%L(\x,\Vlambda) = f(\x) + \sum_i \lambda_i g_i(\x) \leq f(\x).
%\end{myequation*}
%If we now maximize $L(\x,\Vlambda)$ with respect to $\Vlambda$, the best strategy is to avoid penalizing the objective at all, which is achieved by setting $\lambda_i = 0$ for all $i$. Hence, for feasible $\x$,
%\begin{myequation*}
%\max_{\Vlambda:\,\lambda_i \geq 0} L(\x,\Vlambda) = f(\x).
%\end{myequation*}
%This shows that, on feasible points, the Lagrangian reproduces exactly the original objective function.
%
%\medskip
%
%The situation changes dramatically when $\x$ is \textbf{infeasible}. Suppose that at least one constraint is violated, so that there exists an index $i$ such that
%\begin{myequation*}
%g_i(\x) > 0.
%\end{myequation*}
%Since the corresponding multiplier $\lambda_i$ is constrained only to be non-negative, we can increase it arbitrarily. As $\lambda_i \to +\infty$, the penalty term $\lambda_i g_i(\x)$ grows without bound, and therefore the Lagrangian diverges:
%\begin{myequation*}
%\max_{\Vlambda:\,\lambda_i \geq 0} L(\x,\Vlambda) = +\infty.
%\end{myequation*}
%Thus, infeasible points are assigned an infinite cost once the maximization over multipliers is taken into account.
%
%\medskip
%
%Combining these two observations leads to a key insight. The operation
%\begin{myequation*}
%\max_{\Vlambda:\,\lambda_i \geq 0} L(\x,\Vlambda)
%\end{myequation*}
%acts as a \emph{filter}: it leaves the objective value unchanged on feasible points and discards infeasible ones by assigning them infinite cost. Consequently, minimizing this quantity over $\x$ is equivalent to minimizing $f(\x)$ over the feasible set:
%\begin{myequation*}
%\min_{\x} \max_{\Vlambda:\,\lambda_i \geq 0} L(\x,\Vlambda)
%= \min_{\x \in \Omega} f(\x)
%= P^*.
%\end{myequation*}
%This expression recovers the primal problem from the Lagrangian formulation.
%
%\medskip
%
%The dual problem arises when the order of minimization and maximization is reversed. Instead of first enforcing feasibility via the maximization over $\Vlambda$, we first minimize the Lagrangian over $\x$ for a fixed choice of multipliers. For any $\Vlambda \succeq 0$, this produces a value
%\begin{myequation*}
%\min_{\x} L(\x,\Vlambda),
%\end{myequation*}
%which, by construction, cannot exceed the primal optimum. Maximizing this quantity over $\Vlambda$ therefore yields the tightest possible lower bound obtainable through the Lagrangian:
%\begin{myequation*}
%\max_{\Vlambda:\,\lambda_i \geq 0} \min_{\x} L(\x,\Vlambda) \leq P^*.
%\end{myequation*}
%This inequality is precisely the statement of \alert{weak duality}.
%
%
%\medskip
%
%\subsubsection*{Strong Duality}


%Weak duality tells us that the dual problem always provides a lower bound on the primal optimum. In general, however, this bound may be strict: the best value attainable by the dual problem may remain strictly smaller than the optimal value of the primal problem. The difference
%\begin{myequation*}
%p^* - d^*
%\end{myequation*}
%is called the \emph{duality gap} and quantifies how far the two formulations are from each other.
%
%\medskip

A remarkable and highly nontrivial fact is that, 
under suitable conditions --- most notably when $f$ and the $g_i$ are convex and a strictly
feasible point exists (Slater's condition) --- the duality gap closes and
\begin{myequation*}
d^* \;=\; p^*.
\end{myequation*}
This is \alert{strong duality}. When it holds, solving the dual problem is equivalent to
solving the primal: one can recover the primal optimal $\x^*$ from the dual optimal
$\Vlambda^*$. The hard-margin SVM objective is a convex quadratic program with linear
constraints, so Slater's condition is satisfied whenever the data are linearly separable, and
strong duality holds. This justifies the derivation of the dual SVM formulation, to which we
turn next.

\medskip
From a practical standpoint, strong duality has profound consequences. Since the dual problem often has a simpler structure than the primal—fewer variables, simpler constraints, or better numerical properties—it can be advantageous to solve the dual problem instead. When strong duality holds, solving the dual is not merely an approximation: it yields the exact optimal value of the primal problem. 

\medskip

\subsection*{Karush--Kuhn--Tucker (KKT) conditions}

The Karush--Kuhn--Tucker (KKT) conditions provide a precise and powerful characterization of optimal solutions to constrained optimization problems. Conceptually, they can be seen as the natural generalization of the familiar condition ``the gradient must vanish at an optimum'' from unconstrained optimization to the case in which constraints restrict the feasible region.

\medskip

For a convex optimization problem with inequality constraints $g_i(\x)\le 0$, the KKT conditions are \emph{necessary and sufficient} for global optimality. That is, a point $(\x^*,\Vlambda^*)$ is primal--dual optimal if and only if it satisfies the following conditions.

\begin{enumerate}
\item \textbf{Stationarity (first-order optimality).}

The gradient of the Lagrangian with respect to the primal variables must vanish at the optimum:
\begin{myequation*}
\nabla_{\x} L(\x, \Vlambda)\big|_{\x^*,\,\Vlambda^*}
= \nabla f(\x)\big|_{\x^*} + \sum_{i=1}^k \lambda_i^* \nabla g_i(\x)\big|_{\x^*}
= \mathbf{0}.
\end{myequation*}
This condition generalizes the unconstrained requirement $\nabla f(\x)\big|_{\x^*}=\mathbf{0}$. The additional terms $\lambda_i^*\nabla g_i(\x)\big|_{\x^*}$ represent the forces exerted by the active constraints. At the optimum, these forces exactly balance the gradient of the objective, preventing further descent in feasible directions.

\item \textbf{Dual feasibility.}

The Lagrange multipliers must be nonnegative:
\begin{myequation*}
\lambda_i^* \ge 0 \qquad i=1,\ldots,k.
\end{myequation*}
This reflects the fact that constraint violations can only be penalized, never rewarded. A negative multiplier would encourage violating a constraint, which is incompatible with the original constrained problem.

\item \textbf{Primal feasibility.}

The optimal primal solution must satisfy all constraints:
\begin{myequation*}
g_i(\x^*) \le 0 \qquad i=1,\ldots,k.
\end{myequation*}
This is simply the requirement that $\x^*$ belongs to the feasible set of the original problem.

\item \textbf{Complementary slackness.}

For each constraint, the product of the multiplier and the constraint value must vanish:
\begin{myequation*}
\lambda_i^*\, g_i(\x^*) = 0 \qquad i=1,\ldots,k.
\end{myequation*}
\end{enumerate}

\medskip

\textbf{Interpretation of complementary slackness.}

Complementary slackness provides one of the most intuitive and informative insights into constrained optimization. It states that, at the optimum, each constraint must be in exactly one of two regimes:
\begin{itemize}
\item the constraint is \alert{active}, meaning $g_i(\x^*)=0$, in which case the associated multiplier $\lambda_i^*$ may be strictly positive and contributes to the stationarity condition;
\item the constraint is \alert{inactive}, meaning $g_i(\x^*)<0$, in which case the multiplier must vanish, $\lambda_i^*=0$, and the constraint plays no role in determining the optimum.
\end{itemize}
In intuitive terms, only constraints that are ``tight'' at the solution influence the optimal point. Constraints that are strictly satisfied are irrelevant locally: removing them would not change the solution.

\medskip

The KKT conditions above characterize stationary points of the constrained problem. To ensure that such a point is indeed a minimum (and not, for example, a saddle point), a second-order condition is required. Specifically, the Hessian of the Lagrangian with respect to the primal variables,
\begin{myequation*}
H_{\x} = \nabla_{\x}^2 L(\x^*,\Vlambda^*),
\end{myequation*}
must be positive semidefinite on the space of feasible directions. 

For convex optimization problems, this requirement is however automatically satisfied: the objective function is convex, the constraint functions are convex, and therefore the Lagrangian is convex in $\x$. As a result, any point satisfying the KKT conditions is guaranteed to be a global minimum.


\medskip

\subsubsection*{The Dual Problem in Standard Form}

Given the KKT conditions, the dual problem can be expressed explicitly. Define the dual function:
\begin{myequation*}
D(\Vlambda) = \min_{\x} L(\x, \Vlambda)
\end{myequation*}

The dual problem is:
\begin{myalign}
&\max_{\Vlambda} \, D(\Vlambda)\\
&\text{subject to} \quad \lambda_i \geq 0 \quad \text{for } i = 1, \ldots, k
\end{myalign}

The dual problem is always concave in $\Vlambda$ (even if the primal is not convex), which makes it easier to solve in some cases. Combined with strong duality, solving the dual can yield the primal optimum.

\medskip

This formulation is not an ad hoc construction, but follows directly and systematically from the considerations developed in the previous sections. In particular, it is the natural consequence of three key observations:
\begin{itemize}
\item the Lagrangian provides, for each fixed $\Vlambda \succeq 0$, a lower bound on the primal objective;
\item minimizing the Lagrangian over $\x$ extracts the \emph{best} such bound associated with those multipliers;
\item maximizing over $\Vlambda$ selects the tightest bound among all admissible penalties.
\end{itemize}



%\subsection*{Practical Significance for SVMs}
%
%The duality framework is especially valuable for Support Vector Machines, both from a theoretical and a practical standpoint:
%
%\begin{itemize}
%\item The primal SVM problem optimizes the weight vector $\w$ (and bias $b$) in feature space, whose dimension may be very large. By contrast, the dual problem involves only one variable $\lambda_i$ per training example, shifting the complexity from the feature dimension to the number of data points.
%
%\item The dual formulation expresses the optimization problem entirely in terms of inner products between training examples. This structural feature is what enables the use of \alert{kernel methods}, allowing SVMs to operate implicitly in high-dimensional or even infinite-dimensional feature spaces.
%
%\item Complementary slackness reveals that only a subset of training points have nonzero optimal multipliers $\lambda_i^*$. These points, called \alert{support vectors}, are the only ones that influence the final classifier. This sparsity improves computational efficiency and enhances interpretability.
%
%\item The structure of the dual problem admits specialized and highly efficient optimization algorithms, such as Sequential Minimal Optimization (SMO), which are well suited for large-scale learning problems.
%\end{itemize}

\subsection*{Application to SVM: The Dual Problem}

This section derives the dual formulation of the hard-margin SVM optimization problem in a rigorous manner, starting from the primal problem and applying the general duality and KKT framework developed earlier. The derivation highlights how the structure of the dual naturally emerges from the constrained formulation and why it is particularly well suited for kernel-based learning.

\subsubsection*{Primal Problem Setup}

Recall the primal SVM optimization problem:
\begin{myalign}
&\min_{\w,b} \ \frac{1}{2}\|\w\|^2\\
&\text{subject to} \quad t_i(\w^T\x_i + b) \geq 1 \quad \text{for } i=1,\ldots,n.
\end{myalign}

This is a convex optimization problem with linear inequality constraints. In the notation of the general framework:
\begin{itemize}
\item The objective function is
\begin{myequation*}
f(\w,b) = \frac{1}{2}\|\w\|^2,
\end{myequation*}
which is strictly convex in $\w$.
\item The constraint functions are
\begin{myequation*}
g_i(\w,b) = -\big(t_i(\w^T\x_i + b) - 1\big) \le 0,
\quad i=1,\ldots,n,
\end{myequation*}
\item The feasible region is a convex polyhedron defined by these linear constraints.
\end{itemize}

Since the problem is convex and Slater’s condition is satisfied whenever the data are linearly separable, strong duality holds.


\subsubsection*{Formulating the Lagrangian}

Introducing a nonnegative Lagrange multiplier $\lambda_i \ge 0$ for each constraint, the Lagrangian is
\begin{myalign}
L(\w,b,\Vlambda)
&= \frac{1}{2}\|\w\|^2
+ \sum_{i=1}^{n} \lambda_i \, g_i(\w,b)\\
&= \frac{1}{2}\|\w\|^2
- \sum_{i=1}^{n} \lambda_i \big(t_i(\w^T\x_i + b) - 1\big).
\end{myalign}

Expanding the terms yields
\begin{myalign}
L(\w,b,\Vlambda)
= \frac{1}{2}\w^T\w
- \sum_{i=1}^{n}\lambda_i t_i \w^T\x_i
- b\sum_{i=1}^{n}\lambda_i t_i
+ \sum_{i=1}^{n}\lambda_i,
\end{myalign}
with $\Vlambda = (\lambda_1,\ldots,\lambda_n) \succeq 0$.

\medskip

%\subsubsection*{Deriving the Dual via Stationarity}
%
%Recall that the dual function is defined as the minimum of the Lagrangian over the primal
%variables for each fixed choice of multipliers:
%\begin{myequation*}
%q(\Vlambda) \;=\; \min_{\w, b}\, L(\w, b, \Vlambda).
%\end{myequation*}
%This minimization is not an additional modelling choice but is literally what the dual
%function is: for each $\Vlambda \ge 0$, $q(\Vlambda)$ is the lowest value the Lagrangian
%can take when $(\w, b)$ are chosen freely, without any feasibility constraint. As
%established by weak duality, this quantity provides a lower bound on the primal optimum
%$p^*$ for every $\Vlambda \ge 0$, and the dual problem seeks the tightest such bound by
%maximizing $q(\Vlambda)$ over $\Vlambda \ge 0$.
%
%To maximize $q$, we first need to be able to evaluate it, which requires carrying out the
%inner minimization over $(\w, b)$ explicitly. Because the SVM Lagrangian is convex and
%differentiable in $(\w, b)$, this unconstrained minimization has a unique solution,
%characterized by the \alert{stationarity conditions}: the optimal $(\w, b)$ is precisely
%the point at which the gradient of the Lagrangian with respect to the primal variables
%vanishes,
%\begin{myequation*}
%\nabla_{\w} L(\w, b, \Vlambda) = 0, \qquad \frac{\partial L}{\partial b} = 0.
%\end{myequation*}
%These conditions yield explicit expressions for the optimal $(\w^*, b^*)$ as functions of
%the multipliers $\Vlambda$. Substituting these expressions back into the Lagrangian
%eliminates $(\w, b)$ entirely, and what remains is $q(\Vlambda)$ expressed purely in terms
%of the multipliers. The dual problem --- maximizing $q(\Vlambda)$ over $\Vlambda \ge 0$ ---
%then becomes an optimization problem in $\Vlambda$ alone, with no reference to the original
%primal variables.
%
%This substitution is the key step: it converts the minimax problem over both $(\w, b)$ and
%$\Vlambda$ into a simpler maximization over $\Vlambda$ only, and it is what makes the dual
%formulation useful in practice. We now carry out this derivation explicitly.
%\paragraph{Stationarity with respect to $\w$.}
%
%Setting the gradient of $L$ with respect to $\w$ equal to zero gives
%\begin{myequation*}
%\nabla_{\w} L
%= \w - \sum_{i=1}^{n}\lambda_i t_i \x_i
%= \mathbf{0},
%\end{myequation*}
%which implies
%\begin{myequation*}
%\w^* = \sum_{i=1}^{n}\lambda_i t_i \x_i.
%\end{myequation*}
%
%\paragraph{Stationarity with respect to $b$.}
%
%Taking the derivative with respect to $b$ yields
%\begin{myequation*}
%\frac{\partial L}{\partial b}
%= -\sum_{i=1}^{n}\lambda_i t_i
%= 0,
%\end{myequation*}
%which leads to the equality constraint
%\begin{myequation*}
%\sum_{i=1}^{n}\lambda_i t_i = 0.
%\end{myequation*}
%
%%\medskip
%%
%%\textbf{Key insight.}  
%%At optimality, the weight vector $\w^*$ lies in the span of the training examples. This representation, derived directly from the KKT stationarity conditions, is the mathematical foundation of the support vector expansion.
%
%\medskip
%
%\subsubsection*{Explicit Form of the Dual Function}
%
%Substituting the optimal expressions for $\w^*$ and the constraint on $\Vlambda$ into the Lagrangian eliminates the primal variables. We obtain
%\begin{myalign}
%D(\Vlambda)
%&= L(\w^*,b^*,\Vlambda)\\
%&= \frac{1}{2}
%\left(\sum_{i=1}^{n}\lambda_i t_i \x_i\right)^T
%\left(\sum_{j=1}^{n}\lambda_j t_j \x_j\right)
%- \sum_{i=1}^{n}\lambda_i t_i
%\left(\sum_{j=1}^{n}\lambda_j t_j \x_j^T \x_i\right)
%+ \sum_{i=1}^{n}\lambda_i\\
%&= \sum_{i=1}^{n}\lambda_i
%- \frac{1}{2}\sum_{i=1}^{n}\sum_{j=1}^{n}
%\lambda_i\lambda_j t_i t_j \x_i^T\x_j.
%\end{myalign}
%
%This expression depends only on the dual variables $\Vlambda$ and the training data.
%
%\medskip
%
%\subsubsection*{Standard Dual SVM Problem}
%
%The dual SVM problem is therefore:
%\begin{myalign}
%&\max_{\Vlambda}
%\left(
%\sum_{i=1}^{n}\lambda_i
%- \frac{1}{2}\sum_{i=1}^{n}\sum_{j=1}^{n}
%\lambda_i\lambda_j t_i t_j \x_i^T\x_j
%\right)\\
%&\text{subject to} \quad
%\lambda_i \ge 0 \quad \text{for } i=1,\ldots,n,\\
%&\hphantom{\text{subject to}}\quad
%\sum_{i=1}^{n}\lambda_i t_i = 0.
%\end{myalign}
%
%By strong duality, the optimal value of this problem coincides with that of the primal SVM problem.
%
%\medskip
%
%\textbf{Comparison with the primal formulation.}  
%The primal problem involves $d+1$ variables and $n$ constraints, whereas the dual problem involves $n$ variables, one equality constraint, and simple nonnegativity constraints. This shift in structure is central to the effectiveness of SVMs in high-dimensional settings.
%
%\medskip
%
%\subsubsection*{Interpretation of the Dual Variables}
%
%The KKT complementary slackness condition,
%\begin{myequation*}
%\lambda_i^* \big(t_i(\w^{*T}\x_i + b^*) - 1\big) = 0,
%\end{myequation*}
%implies that:
%\begin{itemize}
%\item if $\lambda_i^* > 0$, then $t_i(\w^{*T}\x_i + b^*) = 1$ and the point lies exactly on the margin;
%\item if $t_i(\w^{*T}\x_i + b^*) > 1$, then $\lambda_i^* = 0$ and the point does not influence the classifier.
%\end{itemize}
%
%Thus, only support vectors contribute to the decision function.


\subsubsection*{Deriving the Dual via KKT conditions}

Recall that the dual function is defined as the minimum of the Lagrangian over the primal
variables for each fixed choice of multipliers:
\begin{myequation*}
q(\Vlambda) \;=\; \min_{\w, b}\, L(\w, b, \Vlambda).
\end{myequation*}
This minimization is not an additional modelling choice but is literally what the dual
function is: for each $\Vlambda \ge 0$, $q(\Vlambda)$ is the lowest value the Lagrangian
can take when $(\w, b)$ are chosen freely, without any feasibility constraint. As
established by weak duality, this quantity provides a lower bound on the primal optimum
$p^*$ for every $\Vlambda \ge 0$, and the dual problem seeks the tightest such bound by
maximizing $q(\Vlambda)$ over $\Vlambda \ge 0$.

To maximize $q$, we first need to evaluate it explicitly, which requires carrying out the
inner minimization over $(\w, b)$. The full set of optimality conditions for this
constrained convex program are the \alert{Karush--Kuhn--Tucker (KKT) conditions}, which for
the SVM take the following form:

\begin{enumerate}
\item \textbf{Stationarity:} the gradient of the Lagrangian with respect to the primal
variables vanishes at the optimum,
\begin{myequation*}
\nabla_{\w} L = \mathbf{0}, \qquad \frac{\partial L}{\partial b} = 0;
\end{myequation*}

\item \textbf{Primal feasibility:} the original constraints are satisfied,
\begin{myequation*}
t_i(\w^T \x_i + b) - 1 \;\ge\; 0 \quad \forall\, i;
\end{myequation*}

\item \textbf{Dual feasibility:} the multipliers are non-negative,
\begin{myequation*}
\lambda_i \;\ge\; 0 \quad \forall\, i;
\end{myequation*}

\item \textbf{Complementary slackness:} for each constraint, either the multiplier is zero
or the constraint is active,
\begin{myequation*}
\lambda_i \bigl(t_i(\w^T \x_i + b) - 1\bigr) \;=\; 0 \quad \forall\, i.
\end{myequation*}
\end{enumerate}

These four conditions play distinct roles in the derivation and should not be treated
symmetrically.
\begin{description}
\item[Stationarity] is the condition used actively to derive the dual
function. Because the Lagrangian is convex and differentiable in $(\w, b)$, the inner
minimization over the primal variables is characterized exactly by the vanishing of the
gradient. This yields explicit expressions for the optimal $(\w^*, b^*)$ as functions of
$\Vlambda$, which are then substituted back into the Lagrangian to eliminate the primal
variables and obtain $q(\Vlambda)$ in closed form.

\item[Primal feasibility] does not contribute to the algebraic derivation
of the dual function, but plays an essential role at a higher level: together with the
existence of a strictly feasible point --- which holds whenever the training data are
linearly separable --- it constitutes Slater's condition, which guarantees that strong
duality holds, i.e., $d^* = p^*$. Without this guarantee, the dual optimum would only be a
lower bound on the primal optimum rather than an equal value, and solving the dual would not
be equivalent to solving the primal.

\item[Dual feasibility] is not derived but assumed from the outset. As
established when introducing the Lagrangian, $\lambda_i \ge 0$ is a prerequisite for the
construction to be valid: it ensures that violated constraints are penalized rather than
rewarded. It appears as a constraint in the dual problem from the start.

\item[Complementary slackness] is not used during the derivation of the
dual but becomes essential afterwards, once the dual has been solved: it identifies which
training points are support vectors and is used to recover the optimal bias $b^*$ from the
dual solution.
\end{description}
We now carry out the stationarity-based derivation explicitly.

\paragraph{Stationarity with respect to $\w$.}
The Lagrangian of the hard-margin SVM is
\begin{myequation*}
L(\w, b, \Vlambda)
\;=\; \frac{1}{2}\|\w\|^2
\;-\; \sum_{i=1}^n \lambda_i \bigl[t_i(\w^T \x_i + b) - 1\bigr],
\end{myequation*}
where the constraints have been written as $g_i(\w,b) = 1 - t_i(\w^T\x_i + b) \le 0$,
consistent with the standard form $L = f + \sum_i \lambda_i g_i$.
Setting the gradient with respect to $\w$ to zero gives
\begin{myequation*}
\nabla_{\w} L \;=\; \w - \sum_{i=1}^{n}\lambda_i t_i \x_i \;=\; \mathbf{0},
\end{myequation*}
which yields the first stationarity relation:
\begin{myequation}
\label{eq:w_star}
\w^* \;=\; \sum_{i=1}^{n}\lambda_i t_i \x_i.
\end{myequation}
This is a fundamental structural result: the optimal weight vector is a linear combination
of the training points, with coefficients $\lambda_i t_i$. Points with $\lambda_i = 0$
contribute nothing to $\w^*$. As we will see from complementary slackness, these are
precisely the non-support vectors.

\paragraph{Stationarity with respect to $b$.}
Taking the derivative with respect to $b$ and setting it to zero gives
\begin{myequation*}
\frac{\partial L}{\partial b} \;=\; -\sum_{i=1}^{n}\lambda_i t_i \;=\; 0,
\end{myequation*}
which imposes the equality constraint
\begin{myequation}
\label{eq:sum_constraint}
\sum_{i=1}^{n}\lambda_i t_i \;=\; 0.
\end{myequation}
This condition states that the multipliers, weighted by the class labels, must balance
across the two classes. It will appear as a linear equality constraint in the dual problem.

\subsubsection*{Explicit Form of the Dual Function}

We now substitute the expressions obtained back into the Lagrangian
to eliminate the primal variables entirely. Replacing $\w$ with $\w^*$:
\begin{myalign}
L(\w^*, b, \Vlambda)
&\;=\; \frac{1}{2}\|\w^*\|^2
- \sum_{i=1}^n \lambda_i t_i (\w^{*T}\x_i + b) + \sum_{i=1}^n \lambda_i \notag\\
&\;=\; \frac{1}{2}\sum_{i,j}\lambda_i\lambda_j t_i t_j \x_i^T\x_j
- \sum_{i,j}\lambda_i\lambda_j t_i t_j \x_i^T\x_j
- b\sum_{i=1}^n \lambda_i t_i
+ \sum_{i=1}^n \lambda_i.
\end{myalign}
The term involving $b$ vanishes and the remaining quadratic
terms combine to give the dual function:
\begin{myequation*}
D(\Vlambda) \;=\; \sum_{i=1}^{n}\lambda_i
\;-\; \frac{1}{2}\sum_{i=1}^{n}\sum_{j=1}^{n}
\lambda_i\lambda_j t_i t_j\, \x_i^T\x_j.
\end{myequation*}
This expression depends only on the dual variables $\Vlambda$ and on the training data
through the inner products $\x_i^T \x_j$. The primal variables $(\w, b)$ no longer appear.

\subsubsection*{The Dual SVM Problem}

Maximizing the dual function subject to the constraints derived above gives the dual SVM
problem:
\begin{myalign}
&\max_{\Vlambda} \;\left(
\sum_{i=1}^{n}\lambda_i
- \frac{1}{2}\sum_{i=1}^{n}\sum_{j=1}^{n}
\lambda_i\lambda_j t_i t_j\, \x_i^T\x_j
\right)\\
&\text{subject to} \quad \lambda_i \ge 0 \quad \forall\, i,\\
&\hphantom{\text{subject to}} \quad \sum_{i=1}^{n}\lambda_i t_i = 0.
\end{myalign}
The non-negativity constraints come from dual feasibility, which was required from the
outset. The equality constraint comes from the stationarity condition with respect to $b$.


The structural shift from primal to dual is significant. The primal has $d+1$ variables and
$n$ constraints; the dual has $n$ variables, one equality constraint, and simple
non-negativity constraints. When $d \gg n$ or $d = \infty$, the dual is far more tractable.
Moreover, the dual depends on the data only through inner products $\x_i^T \x_j$, which is
what enables the kernel trick.

\subsubsection*{Recovering the Solution via Complementary Slackness}

Once the dual problem has been solved and the optimal $\Vlambda^*$ obtained, the primal
solution is recovered in two steps. By strong duality --- guaranteed by primal feasibility and Slater's condition --- the
optimal value of this problem equals $p^*$:  the primal solution is then recovered
via the first stationarity relation and T
the weight vector follows immediately as
$\w^* \;=\; \sum_{i=1}^{n}\lambda_i^{*} t_i \x_i$. 

\medskip
The bias $b^*$ is recovered using the complementary slackness
condition,
\begin{myequation*}
\lambda_i^* \bigl(t_i(\w^{*T}\x_i + b^*) - 1\bigr) \;=\; 0 \quad \forall\, i,
\end{myequation*}
which states that for any support vector $\x_i$ with $\lambda_i^* > 0$, the constraint must
be active: $t_i(\w^{*T}\x_i + b^*) = 1$. 

Solving for $b^*$ yields
\begin{myalign}
b^* &= \frac{1}{t_k} - \w^{*T}\x_k
= t_k - \w^{*T}\x_k \qquad (\text{since } t_k^2 = 1)\\
&= t_k - \sum_{j=1}^{n} \lambda_j^* t_j \x_j^T \x_k.
\end{myalign}

In kernel form, this becomes
\begin{myequation*}
b^* = t_k - \sum_{j=1}^{n} \lambda_j^* t_j \kappa(\x_j, \x_k).
\end{myequation*}

Because multiple support vectors may exist, numerical stability is improved by averaging this expression over all them:
\begin{myequation*}
b^* =
\frac{1}{|\mathcal{S}|}
\sum_{k \in \mathcal{S}}
\left(
t_k - \sum_{j \in \mathcal{S}} \lambda_j^* t_j \kappa(\x_j, \x_k)
\right),
\end{myequation*}
where
\begin{myequation*}
\mathcal{S} = \{\, i \mid \lambda_i^* > 0 \,\}
\end{myequation*}
denotes the set of support vector indices.



%Complementary slackness also provides the full interpretation of the dual solution. For
%each training point, exactly one of two cases holds:
%
%\begin{itemize}
%\item \textbf{$\lambda_i^* = 0$:} the point does not contribute to $\w^*$ and lies
%strictly outside the margin, $t_i(\w^{*T}\x_i + b^*) > 1$. It plays no role in defining
%the decision boundary and could be removed without changing the solution.
%
%\item \textbf{$\lambda_i^* > 0$:} the constraint is active and the point lies exactly on
%the margin, $t_i(\w^{*T}\x_i + b^*) = 1$. It is a \alert{support vector} and contributes
%to $\w^*$ with weight $\lambda_i^* t_i$.
%\end{itemize}
%
%The sparsity of the SVM solution --- the fact that the classifier is determined by a
%typically small subset of training points --- is therefore not a computational convenience
%but a direct and necessary consequence of the KKT complementary slackness condition.

\subsubsection*{Kernel Formulation}

A decisive advantage of the dual formulation is that the data appear only through inner products $\x_i^T\x_j$. This allows us to replace them with a kernel function
\begin{myequation*}
\kappa(\x_i,\x_j) = \x_i^T\x_j,
\end{myequation*}
or, more generally,
\begin{myequation*}
\kappa(\x_i,\x_j)
= \boldsymbol{\phi}(\x_i)^T \boldsymbol{\phi}(\x_j),
\end{myequation*}
where $\boldsymbol{\phi}$ is an implicit feature map.

The dual problem in kernelized form becomes:
\begin{myalign}
&\max_{\Vlambda}
\left(
\sum_{i=1}^{n}\lambda_i
- \frac{1}{2}\sum_{i=1}^{n}\sum_{j=1}^{n}
\lambda_i\lambda_j t_i t_j \kappa(\x_i,\x_j)
\right)\\
&\text{subject to} \quad
\lambda_i \ge 0 \quad \text{for } i=1,\ldots,n,\\
&\hphantom{\text{subject to}}\quad
\sum_{i=1}^{n}\lambda_i t_i = 0.
\end{myalign}

This formulation enables SVMs to learn nonlinear decision boundaries while preserving the convexity and strong duality properties of the original optimization problem.


\medskip

\subsubsection*{Support Vectors and Complementary Slackness}

By the complementary slackness condition, for each training example we have
\begin{myequation*}
\lambda_i^* \big(t_i(\w^{*T}\x_i + b^*) - 1\big) = 0.
\end{myequation*}

This condition establishes a precise correspondence between the optimal dual variables and the geometric role of the training points:

\begin{itemize}
\item If $\lambda_i^* > 0$, then necessarily
\begin{myequation*}
t_i(\w^{*T}\x_i + b^*) = 1,
\end{myequation*}
meaning that the point $\x_i$ lies exactly on the margin boundary. Such points are called \alert{support vectors}.

\item If $\lambda_i^* = 0$, then the constraint is inactive and
\begin{myequation*}
t_i(\w^{*T}\x_i + b^*) \ge 1.
\end{myequation*}
The point is correctly classified with a margin strictly larger than one and does not influence the optimal separating hyperplane.
\end{itemize}

Thus, complementary slackness provides a rigorous explanation for the sparsity of the SVM solution: only training points that lie on the margin contribute to defining the classifier. In practice, the number of support vectors is typically much smaller than the total number of training examples $n$, yielding compact models and efficient prediction.


\subsection*{Classification Using the Trained SVM}

Once the optimal parameters $\w^*$ and $b^*$ have been determined, a new point $\x$ is classified using the decision function
\begin{myequation*}
y(\x) = \w^{*T}\x + b^*.
\end{myequation*}

Substituting the support vector expansion of $\w^*$ gives
\begin{myequation*}
y(\x) = \sum_{i=1}^{n} \lambda_i^* t_i \x_i^T \x + b^*,
\end{myequation*}
and the predicted label is $\mathrm{sign}(y(\x))$.

\medskip

Using a kernel function, the decision function becomes
\begin{myequation*}
y(\x) = \sum_{i=1}^{n} \lambda_i^* t_i \kappa(\x_i, \x) + b^*.
\end{myequation*}

By sparsity of the solution, this expression simplifies to
\begin{myequation*}
y(\x) = \sum_{i \in \mathcal{S}} \lambda_i^* t_i \kappa(\x_i, \x) + b^*,
\end{myequation*}
where the sum runs only over support vectors.

\medskip

\noindent
\textbf{Computational efficiency.}
The decision function requires kernel evaluations only between the new point and the support vectors, rather than the entire training set. When $|\mathcal{S}| \ll n$, prediction is therefore fast and memory-efficient.

\medskip

\noindent
\textbf{Key advantages of this formulation:}
\begin{itemize}
\item The classifier depends on the data exclusively through kernel evaluations, enabling nonlinear decision boundaries without explicit feature construction.
\item The sparsity induced by complementary slackness yields interpretable models and efficient computation at both training and test time.
\item The kernel-based dual formulation illustrates how general principles of convex duality translate directly into practical and powerful learning algorithms.
\end{itemize}

\subsection*{Soft-Margin SVM: Handling Non-Separable Data}

The hard-margin SVM formulation developed in the previous section relies on the assumption that the training data are linearly separable in feature space. While this assumption is mathematically convenient and leads to a clean geometric picture, it is rarely satisfied in practice. Real-world data are often noisy, may contain outliers, or may simply not admit a perfect separating hyperplane, even after mapping into a high-dimensional feature space via a kernel.

The soft-margin SVM extends the hard-margin framework in a principled way to cope with such situations. The key idea is to relax the strict separation constraints, allowing controlled violations, while still preserving the core philosophy of margin maximization.

\medskip

Recall that the hard-margin SVM enforces the constraints
\begin{myequation*}
t_i(\w^T\boldsymbol{\phi}(\x_i)+b) \geq 1 \quad \text{for all } i=1,\ldots,n.
\end{myequation*}

Geometrically, these inequalities require every training point to lie on the correct side of the separating hyperplane and at least at unit distance (in the scaled sense) from it. When the data are not linearly separable, this system of inequalities becomes infeasible: no choice of $(\w,b)$ can satisfy all constraints simultaneously. As a consequence, the hard-margin optimization problem admits no solution and completely breaks down.

From a modeling perspective, this behavior is undesirable. A single mislabeled point or a small amount of noise should not prevent the learning algorithm from producing a reasonable classifier. This motivates a relaxation of the constraints that tolerates violations, but does so in a controlled and quantifiable way.

\medskip

\subsubsection*{Slack Variables and Relaxed Constraints}

Constraint relaxation is introduced through the use of \alert{slack variables}. For each training example, we associate a nonnegative slack variable $\xi_i \geq 0$ that measures the degree to which the margin constraint is violated. The hard constraint is replaced by
\begin{myequation*}
t_i(\w^T\boldsymbol{\phi}(\x_i)+b) \geq 1 - \xi_i \quad \text{for all } i=1,\ldots,n.
\end{myequation*}

The role of $\xi_i$ can be understood geometrically by examining different regimes. If $\xi_i = 0$, the constraint reduces to the hard-margin condition and the point is correctly classified with the required margin. If $0 < \xi_i < 1$, the point is still correctly classified, but it lies inside the margin strip and therefore enjoys a reduced margin. When $\xi_i = 1$, the point lies exactly on the decision boundary. Finally, if $\xi_i > 1$, the point is misclassified and lies on the wrong side of the hyperplane.

In this way, the slack variables provide a continuous measure of constraint violation, allowing the optimization problem to remain feasible even when perfect separation is impossible.

\medskip

\subsubsection*{The Soft-Margin Optimization Problem}

Introducing slack variables alone is not sufficient: we must also discourage excessive violations. This is achieved by augmenting the objective function with a penalty term that grows with the total slack. The resulting optimization problem is
\begin{myalign}
&\min_{\w,b,\Vxi} \ \frac{1}{2}\|\w\|^2 + C\sum_{i=1}^{n}\xi_i\\
&\text{subject to} \quad t_i(\w^T\boldsymbol{\phi}(\x_i)+b) \geq 1-\xi_i \quad \text{for } i=1,\ldots,n,\\
&\hphantom{\text{subject to}}\quad \xi_i \geq 0 \quad \text{for } i=1,\ldots,n,
\end{myalign}
where $\Vxi = (\xi_1,\ldots,\xi_n)$ and $C>0$ is a regularization parameter.

The structure of the objective reflects a trade-off between two competing goals. The term $\frac{1}{2}\|\w\|^2$ favors large margins, as in the hard-margin case. The term $C\sum_i \xi_i$ penalizes violations of the margin constraints. The parameter $C$ controls the relative importance of these objectives. For large values of $C$, violations are heavily penalized and the solution prioritizes correct classification of the training data, potentially at the expense of a smaller margin. For small values of $C$, violations are penalized less strongly, leading to wider margins but possibly more misclassifications. In practice, $C$ is treated as a hyperparameter and selected by cross-validation.

\medskip

\subsubsection*{Lagrangian Formulation}

To derive the dual problem and characterize optimal solutions, we construct the Lagrangian associated with the soft-margin primal problem. Two sets of inequality constraints are present: the relaxed margin constraints and the nonnegativity constraints on the slack variables. Accordingly, we introduce Lagrange multipliers $\lambda_i \geq 0$ for the former and $\alpha_i \geq 0$ for the latter. The Lagrangian is
\begin{myalign}
L(\w,b,\Vxi,\Vlambda,\Valpha)
&= \frac{1}{2}\|\w\|^2 + C\sum_{i=1}^{n}\xi_i\\
&\quad - \sum_{i=1}^{n}\lambda_i\left(t_i(\w^T\boldsymbol{\phi}(\x_i)+b) - 1 + \xi_i\right)\\
&\quad - \sum_{i=1}^{n}\alpha_i\xi_i.
\end{myalign}

This expression combines the objective function with weighted penalties for constraint violations, in direct analogy with the hard-margin case, but now enriched by the presence of slack variables.

\medskip

\subsubsection*{KKT Conditions}

Optimality of a primal--dual solution is characterized by the Karush--Kuhn--Tucker conditions. Stationarity of the Lagrangian with respect to the primal variables yields
\begin{myalign}
\nabla L|_{\w} &= \mathbf{0},\\
\frac{\partial L}{\partial b} &= 0,\\
\frac{\partial L}{\partial \xi_i} &= 0 \quad \text{for } i=1,\ldots,n.
\end{myalign}

These are complemented by primal feasibility of the constraints, dual feasibility of the multipliers, and complementary slackness conditions:
\begin{myalign}
t_i(\w^{*T}\boldsymbol{\phi}(\x_i)+b^*) - 1 + \xi_i^* &\geq 0,\\
\xi_i^* &\geq 0,\\
\lambda_i^*,\alpha_i^* &\geq 0,\\
\lambda_i^*\left(t_i(\w^{*T}\boldsymbol{\phi}(\x_i)+b^*) - 1 + \xi_i^*\right) &= 0,\\
\alpha_i^*\xi_i^* &= 0.
\end{myalign}

Together, these conditions encode both feasibility and optimal balance between margin maximization and violation penalties.

\medskip

\subsubsection*{Deriving the Dual Problem}

Evaluating the stationarity conditions explicitly yields
\begin{myalign}
\w^* &= \sum_{i=1}^{n}\lambda_i t_i \boldsymbol{\phi}(\x_i),\\
\sum_{i=1}^{n}\lambda_i t_i &= 0,\\
C - \lambda_i - \alpha_i &= 0.
\end{myalign}

The last relation implies $\lambda_i \leq C$ for all $i$, since $\alpha_i \geq 0$. Substituting the optimal expressions into the Lagrangian and simplifying as in the hard-margin case leads to the dual problem
\begin{myalign}
&\max_{\Vlambda}
\left(
\sum_{i=1}^{n}\lambda_i
- \frac{1}{2}\sum_{i=1}^{n}\sum_{j=1}^{n}
\lambda_i\lambda_j t_i t_j \kappa(\x_i,\x_j)
\right)\\
&\text{subject to} \quad
0 \leq \lambda_i \leq C \quad \text{for } i=1,\ldots,n,\\
&\hphantom{\text{subject to}}\quad
\sum_{i=1}^{n}\lambda_i t_i = 0.
\end{myalign}

The only formal difference with respect to the hard-margin dual problem is the upper bound $\lambda_i \leq C$. In the limit $C \to \infty$, this constraint disappears and the hard-margin formulation is recovered.

\medskip

\subsubsection*{Interpretation via Complementary Slackness}

Complementary slackness provides a detailed picture of the role played by each training example. If $\lambda_i^* = 0$ and $\xi_i^* = 0$, the point lies outside the margin and does not influence the solution. If $0 < \lambda_i^* < C$, then $\xi_i^* = 0$ and the point lies exactly on a margin boundary: these are classical support vectors. If $\lambda_i^* = C$, then $\xi_i^* > 0$ and the point either lies inside the margin or is misclassified. Such points are also support vectors, but they correspond to margin violations.

Thus, as in the hard-margin case, only support vectors contribute to the expansion of the optimal weight vector
\begin{myequation*}
\w^* = \sum_{i \in \mathcal{S}} \lambda_i^* t_i \boldsymbol{\phi}(\x_i),
\end{myequation*}
where $\mathcal{S}$ denotes the set of indices with $\lambda_i^* > 0$.

\medskip

\subsubsection*{Recovering Primal Parameters and Classification}

The bias term $b^*$ can be recovered using any support vector with $0 < \lambda_i^* < C$:
\begin{myequation*}
b^* = t_k - \sum_{j \in \mathcal{S}}\lambda_j^* t_j \kappa(\x_j,\x_k).
\end{myequation*}
In practice, averaging this expression over all such support vectors improves numerical stability.

Classification of a new point $\x$ then proceeds as
\begin{myequation*}
y(\x) = \sum_{i \in \mathcal{S}}\lambda_i^* t_i \kappa(\x_i,\x) + b^*,
\end{myequation*}
with predicted label $\text{sign}(y(\x))$.

\medskip

The soft-margin SVM thus preserves the key structural advantages of the hard-margin formulation—sparsity, kernelization, and convex optimization—while extending its applicability to realistic, noisy, and non-separable datasets.


\subsection*{Extensions of SVM}

So far, we have developed Support Vector Machines as a method for binary classification, first in the hard-margin setting and then in the soft-margin setting. In both cases, the central ideas were the same: represent the classifier as a separating hyperplane in a (possibly high-dimensional) feature space, maximize the margin, and exploit the dual formulation to obtain a sparse kernelized decision function.

In many applications, however, the learning task departs from the basic binary setting. Two particularly important generalizations are (i) classification with more than two classes and (ii) prediction of continuous outputs rather than discrete labels. In what follows, we describe how the SVM philosophy extends to these settings while preserving the key structural advantages of the method.

\medskip

\subsubsection*{Multiclass Classification}

The standard SVM is designed for binary labels $t_i \in \{-1,+1\}$. When the dataset contains $K>2$ classes, there is no unique ``canonical'' generalization of the margin constraint, because the notion of separation must now distinguish among several categories. A common and effective strategy is therefore to reduce the multiclass problem to a collection of binary SVM problems, each of which can be solved with exactly the machinery already developed.

\medskip

\textbf{One-vs-All (OvA) approach.}

The most direct reduction is the \alert{one-vs-all} (or one-vs-rest) scheme. The idea is to build $K$ binary classifiers, one for each class. For a fixed class $k \in \{1,\ldots,K\}$, we relabel the training set by declaring class $k$ as the positive class and all other classes as negative. Denoting these induced binary labels by $t_{i,k}$, we have $t_{i,k}=+1$ if the original label of $\x_i$ is $k$, and $t_{i,k}=-1$ otherwise.

For each class $k$, we then train a standard (soft-margin) SVM classifier with parameters $(\w_k,b_k)$ and slacks $\Vxi_k$. Concretely, we solve
\begin{myalign}
&\min_{\w_k, b_k, \Vxi_k} \frac{1}{2}\|\w_k\|^2 + C\sum_{i=1}^{n}\xi_{i,k}\\
&\text{subject to} \quad t_{i,k}(\w_k^T\boldsymbol{\phi}(\x_i)+b_k) \geq 1-\xi_{i,k} \quad \text{for } i=1,\ldots,n\\
&\quad\quad\quad\quad\xi_{i,k} \geq 0.
\end{myalign}

The output of the $k$-th classifier can be interpreted as a real-valued \emph{score} measuring how confidently $\x$ belongs to class $k$. Using the dual representation, this score takes the familiar kernel form
\begin{myequation*}
y_k(\x) = \sum_{i \in \mathcal{S}_k}\lambda_{i,k}^* t_{i,k} \kappa(\x_i, \x) + b_k^*,
\end{myequation*}
where $\mathcal{S}_k$ denotes the set of support vectors for the $k$-th binary problem.

At prediction time, we evaluate all $K$ scores and assign $\x$ to the class with the largest score:
\begin{myequation*}
\hat{t} = \arg\max_{k \in \{1,\ldots,K\}} y_k(\x).
\end{myequation*}
In words, OvA chooses the class whose ``separate-from-the-rest'' hyperplane considers the point most strongly positive.

\medskip

\textbf{Advantages and limitations.}

The OvA approach is popular because it is simple, modular, and computationally efficient: it requires solving only $K$ binary SVM problems, each of which can be trained with the same algorithms used in the binary case. Moreover, the $K$ classifiers are independent, so one may use different kernels or different regularization parameters $C$ for different classes if the application suggests doing so.

At the same time, the reduction is not without caveats. Since class $k$ typically contains $n_k$ examples, the induced binary dataset for classifier $k$ contains $n_k$ positives and $n-n_k$ negatives; if $n_k \ll n$, the learning problem can be highly imbalanced. In addition, the $K$ hyperplanes are trained independently, so their decision regions are not guaranteed to be globally consistent: it is possible for several classifiers to give similar scores, and (in principle, though rarely in practice) all scores could be negative, making the interpretation of the winning class less geometrically clean than in the binary case. Nevertheless, in many applications OvA performs very well and is often the default choice.

\medskip

\textbf{One-vs-One (OvO) alternative.}

A widely used alternative is \alert{one-vs-one}. Instead of separating one class from all the others, OvO trains a binary SVM for every pair of classes. For $K$ classes, this yields
\begin{myequation*}
\binom{K}{2} = \frac{K(K-1)}{2}
\end{myequation*}
classifiers. Each classifier is trained only on examples from its two corresponding classes, which typically reduces imbalance and can lead to cleaner binary separations. During prediction, each pairwise classifier casts a vote, and the class with the highest number of votes is selected. The price paid is computational: training and evaluating $\binom{K}{2}$ classifiers can be more expensive than OvA, especially when $K$ is large.

\medskip

\subsubsection*{Support Vector Regression}

Classification assigns one of finitely many labels to each input point. In many problems, however, the desired output is a real number: predicting prices, temperatures, risk scores, or physical quantities from sensor measurements are typical examples. Support Vector Regression (SVR) extends the SVM framework to this setting while retaining the same core principles: regularize by controlling $\|\w\|$, build a sparse solution in terms of support vectors, and exploit kernel methods to model nonlinear relations.

\medskip

%\textbf{$\epsilon$-insensitive loss.}

The central modeling choice in SVR is the loss function. Instead of requiring a classification margin, SVR introduces an \alert{$\epsilon$-insensitive tube} around the regression function $\hat{y}(\x)=\w^T\boldsymbol{\phi}(\x)+b$. Deviations smaller than $\epsilon$ are treated as negligible, while larger deviations are penalized linearly. This is formalized by the $\epsilon$-insensitive loss
\begin{myequation*}
L_\epsilon(y_i, \hat{y}_i) =
\begin{cases}
0 & \text{if } |y_i - \hat{y}_i| \leq \epsilon\\
|y_i - \hat{y}_i| - \epsilon & \text{otherwise}.
\end{cases}
\end{myequation*}

Geometrically, one may picture a tube of width $2\epsilon$ around the prediction function. Points falling inside the tube incur no loss; only points outside the tube contribute to the training penalty. This choice makes the method robust to small perturbations and encourages sparsity, because only points outside (or exactly on) the tube become support vectors in the dual representation.

\medskip

\textbf{Primal SVR problem.}

To translate this idea into a constrained optimization problem, we introduce two sets of nonnegative slack variables, $\xi_i \ge 0$ and $\xi_i^* \ge 0$, which measure how much the prediction violates the tube above and below, respectively. The constraints enforce that every point lies within the tube up to slack:
\begin{myalign}
y_i - \w^T\boldsymbol{\phi}(\x_i) - b &\leq \epsilon + \xi_i \quad \text{for } i=1,\ldots,n,\\
\w^T\boldsymbol{\phi}(\x_i) + b - y_i &\leq \epsilon + \xi_i^* \quad \text{for } i=1,\ldots,n.
\end{myalign}

The corresponding objective balances flatness of the regression function (controlled by $\|\w\|$) with the total amount of tube violation:
\begin{myalign}
&\min_{\w, b, \Vxi, \Vxi^*} \frac{1}{2}\|\w\|^2 + C\sum_{i=1}^{n}(\xi_i + \xi_i^*)\\
&\text{subject to} \quad y_i - \w^T\boldsymbol{\phi}(\x_i) - b \leq \epsilon + \xi_i \quad \text{for } i=1,\ldots,n\\
&\quad\quad\quad\quad \w^T\boldsymbol{\phi}(\x_i) + b - y_i \leq \epsilon + \xi_i^* \quad \text{for } i=1,\ldots,n\\
&\quad\quad\quad\quad \xi_i, \xi_i^* \geq 0.
\end{myalign}

Here, $\xi_i$ quantifies violations of the upper constraint (the true value lies above the tube), while $\xi_i^*$ quantifies violations of the lower constraint (the true value lies below the tube). The parameter $C$ plays the same conceptual role as in soft-margin classification: large $C$ forces the regression function to fit the data more tightly (allowing fewer and smaller violations), while smaller $C$ tolerates larger deviations in exchange for a smoother (lower-norm) predictor.

\medskip

In the next step, one can proceed exactly as in the classification case: construct the Lagrangian, apply the KKT conditions, and derive a dual problem expressed only through kernel evaluations. The resulting SVR predictor is again a sparse expansion in terms of a subset of training points, reflecting the fact that only points outside the $\epsilon$-tube influence the final regression function.

\section*{Maximum Margin Classifiers and VC Dimension}

The preference for large margins is not only geometrically intuitive but also deeply rooted
in statistical learning theory. We now show that the margin has a precise and quantitative
consequence for the generalization ability of the classifier, through its connection with the
VC dimension. As established in the treatment of the Fundamental Theorem, the VC dimension
is the exact combinatorial quantity that controls generalization: the complexity term in the
generalization bound grows with $\mathrm{VCdim}(\mathcal{H})$ and decays
with the sample size. We will see that maximizing the margin is equivalent to minimizing an
upper bound on the VC dimension of the effective hypothesis class, and that this is the
fundamental mechanism by which SVMs generalize well even in very high-dimensional or
infinite-dimensional feature spaces.

\subsection*{VC Dimension of Linear Classifiers}

We start from stating (without proving it here) that the class of all linear classifiers in $\mathbb{R}^d$, which assigns labels according to
$h_{\mathbf{w},b}(\mathbf{x}) = \mathrm{sign}(\mathbf{w}^\top \mathbf{x} + b)$, has VC
dimension exactly $d+1$.

%\begin{theorem}
%$\mathrm{VCdim}(\mathcal{H}_{\mathrm{lin}}) = d + 1$.
%\end{theorem}
%
%The upper bound follows from Radon's theorem: any set of $d+2$ points in $\mathbb{R}^d$ can
%be partitioned into two subsets whose convex hulls intersect, making the corresponding
%labeling unrealizable by any hyperplane. The lower bound is witnessed by the $d$ standard
%basis vectors together with the origin, which can be shattered.

The VC dimension therefore scales linearly with the ambient dimension $d$. From the
perspective of the generalization bound, this means the complexity term
grows with $d$, requiring a correspondingly large training set for the bound to be
informative. In a kernel-induced feature space, where $d$ may be infinite, the VC dimension
of $\mathcal{H}_{\mathrm{lin}}$ is also infinite and the bound becomes
vacuous. One might therefore expect SVMs with kernels to generalize poorly without enormous
amounts of data. In practice, however, this is not the case. Resolving this contradiction
requires looking beyond the raw VC dimension of the full linear class and accounting for the
additional structure imposed by the margin.

%\subsection*{The Margin Reduces Effective Capacity}

Recall from the previous section that the geometric margin $\gamma = \min_i \gamma_i$
measures the distance from the decision boundary to the nearest training point, and that the
hard-margin SVM finds the unique hyperplane that maximizes this quantity. We now show that
restricting attention to classifiers that achieve a margin of at least $\gamma$ does not
merely provide geometric robustness --- it provably reduces the VC dimension of the
hypothesis class.

Define the class of margin-$\gamma$ classifiers as
\begin{myequation*}
\mathcal{H}_\gamma \;=\; \bigl\{(\mathbf{w}, b) : \|\mathbf{w}\| = 1,\;
\min_i\, t_i(\mathbf{w}^\top \mathbf{x}_i + b) \ge \gamma \bigr\},
\end{myequation*}
a strict subset of $\mathcal{H}_{\mathrm{lin}}$. The following theorem gives a bound on its
VC dimension that depends not on $d$, but only on the geometry of the data and the margin.

\begin{theorem}[Margin-based VC bound]
\label{thm:margin_vc}
Suppose all training points satisfy $\|\mathbf{x}_i\| \le R$. Then
\begin{myequation}
\label{eq:margin_vc}
\mathrm{VCdim}(\mathcal{H}_\gamma) \;\le\; \left\lfloor \frac{R^2}{\gamma^2} \right\rfloor + 1.
\end{myequation}
\end{theorem}

here $R$ is the radius of the smallest enclosing ball of the dataset, centered at its optimal center
\begin{myequation*}
R = \min_{\mathbf{c}} \max_i \|\mathbf{x}_i - \mathbf{c}\|,
\end{myequation*}
the minimax radius of the data cloud. Note that this quantity is translation-invariant by construction: shifting all points by a constant vector $\mathbf{v}$ shifts the optimal center $\mathbf{c}^*$ by the same amount and leaves $R^*$ unchanged.



%\paragraph{Proof sketch.}
%If $\mathcal{H}_\gamma$ could shatter a set of $m$ points, then by a geometric argument
%involving the inner products $\langle \mathbf{w}, \mathbf{x}_i \rangle$, one can show that
%$2^m$ distinct label configurations must each be separated by at least $\gamma$. The norm
%constraint $\|\mathbf{w}\| = 1$ and the radius bound $\|\mathbf{x}_i\| \le R$ together force
%$m \le (R/\gamma)^2$, giving the result.
\medskip
Three features of this bound deserve emphasis.

\begin{itemize}
\item \textbf{Dimension independence.} The bound does not involve $d$
at all. Even in an infinite-dimensional kernel-induced feature space, the VC dimension of
$\mathcal{H}_\gamma$ remains finite provided $\gamma$ is bounded away from zero relative to
$R$. This is the key property that resolves the apparent paradox of kernel SVMs.

\item \textbf{Monotonicity in margin.} The bound decreases strictly as $\gamma$ increases.
A wider margin corresponds to a hypothesis class of provably lower VC dimension and therefore
lower complexity, which by the Fundamental Theorem translates directly into better
generalization guarantees.

\item \textbf{The ratio $R/\gamma$ as the relevant scale.} What matters is not the margin or
the data radius in isolation, but their ratio. A margin that is large relative to the spread
of the data drives the capacity reduction; the same absolute margin value may correspond to
very different effective complexities depending on the scale of the problem.
\end{itemize}

\subsection*{The SVM Generalization Bound}

%Substituting this bound into the VC generalization bound, we
%obtain that with probability at least $1-\delta$ over a training set of $n$ i.i.d.\ samples,
%the SVM classifier satisfies
%\begin{myequation*}
%\label{eq:svm_gen_bound}
%L_{\mathcal{D}}(h) \;\le\; \hat{L}(h) \;+\;
%O\!\left(\sqrt{\frac{\dfrac{R^2}{\gamma^2}\log n + \log(1/\delta)}{n}}\right).
%\end{myequation*}
Substituting the margin-based VC bound into the general VC
generalization bound requires replacing the VC dimension $d$ with the
upper bound $\lfloor R^2/\gamma^2 \rfloor + 1$. Recall that the VC generalization bound states that,
with probability at least $1-\delta$, every $h \in \mathcal{H}$ satisfies
\begin{myequation*}
L_{\mathcal{D}}(h) \;\le\; \hat{L}(h) \;+\;
O\!\left(\sqrt{\frac{d\log(n/d) + \log(1/\delta)}{n}}\right).
\end{myequation*}
Setting $d_{\mathrm{margin}} = \lfloor R^2/\gamma^2 \rfloor + 1$ and substituting gives
\begin{myequation*}
L_{\mathcal{D}}(h) \;\le\; \hat{L}(h) \;+\;
O\!\left(\sqrt{\frac{d_{\mathrm{margin}}\log(n/d_{\mathrm{margin}}) + \log(1/\delta)}{n}}\right).
\end{myequation*}
Expanding the logarithmic factor, we have
\begin{myequation*}
d_{\mathrm{margin}}\log\!\left(\frac{n}{d_{\mathrm{margin}}}\right)
\;=\;
\left\lfloor\frac{R^2}{\gamma^2}\right\rfloor \log\!\left(\frac{n\gamma^2}{R^2}\right)
\;\lesssim\;
\frac{R^2}{\gamma^2}\log\!\left(\frac{n\gamma^2}{R^2}\right)
\;\le\;
\frac{R^2}{\gamma^2}\log n,
\end{myequation*}
where the last inequality holds because $\gamma^2/R^2 \le 1$ in any non-trivial
configuration, so $\log(n\gamma^2/R^2) \le \log n$. Absorbing the constant factors into the
$O(\cdot)$ notation then yields
\begin{myequation}
\label{eq:svm_gen_bound}
L_{\mathcal{D}}(h) \;\le\; \hat{L}(h) \;+\;
O\!\left(\sqrt{\frac{\dfrac{R^2}{\gamma^2}\log n + \log(1/\delta)}{n}}\right),
\end{myequation}
Two aspects of this bound are worth highlighting explicitly.

First, the dependence on the dimension $d$ has disappeared entirely: the complexity
term is governed solely by the ratio $R/\gamma$ and the sample size $n$. This is a direct
consequence of Theorem~\ref{thm:margin_vc}, which bounds the effective VC dimension in terms
of geometry rather than dimension.

Second, the bound decreases as $\gamma$ increases. Specifically, for a fixed $\delta$ the complexity term scales
as $O\!\left(\dfrac{R}{\gamma}\sqrt{\dfrac{\log n}{n}}\right)$, so doubling the margin reduces the
complexity term by a factor of $2$, independently of the dimensionality of the feature
space. This makes precise the intuition that a larger margin yields a better generalization
guarantee: it is not merely a heuristic preference but a quantitative improvement in the
upper bound on the generalization error.

In the separable case, the hard-margin SVM achieves $\hat{L}(h) = 0$, so the generalization
error is controlled entirely by the complexity term, which depends on $d$ only through $R$
and $\gamma$. This resolves the paradox raised above: even though $\mathcal{H}_{\mathrm{lin}}$
has infinite VC dimension in a kernel-induced space, the maximum-margin classifier belongs to
the much smaller subclass $\mathcal{H}_\gamma$ whose effective VC dimension is at most
$(R/\gamma)^2$. It is this quantity --- not the ambient dimension --- that enters the
generalization bound and governs generalization performance.

\subsection*{Margin Maximization as Capacity Control}

The optimization problem introduced earlier for the hard-margin SVM solves
\begin{myequation*}
\min_{\mathbf{w}, b} \;\frac{1}{2}\|\mathbf{w}\|^2 \quad \text{subject to} \quad
t_i(\mathbf{w}^\top \mathbf{x}_i + b) \ge 1, \quad i=1,\ldots,n,
\end{myequation*}
which is equivalent to maximizing the geometric margin $\gamma = 1/\|\mathbf{w}\|$:
\begin{myequation*}
\max_{\mathbf{w}, b} \;\frac{1}{\|\mathbf{w}\|} \;\;\Longleftrightarrow\;\;
\min_{\mathbf{w}, b} \;\|\mathbf{w}\|^2.
\end{myequation*}
By the  this simultaneously minimizes an upper bound on
$\mathrm{VCdim}(\mathcal{H}_\gamma)$, reducing the complexity term
in~\eqref{eq:svm_gen_bound} without increasing the training error. The SVM therefore
implements, through a purely geometric objective, exactly the capacity control that the
Fundamental Theorem identifies as the key condition for learnability: it restricts the
effective hypothesis class to one of provably low VC dimension, ensuring that low training
error is more likely to translate into low generalization error.

This connection also gives a precise account of the role of the support vectors identified
earlier. Since the margin $\gamma$ is determined entirely by the training points closest to
the decision boundary --- the support vectors --- it is these points alone that control the
effective VC dimension of the learned classifier through the bound~\eqref{eq:margin_vc}. The
sparsity of the SVM solution is therefore not merely a computational convenience: it reflects
the fact that generalization is governed by a small number of critical points that determine
the margin, and hence the capacity of the hypothesis class.

\subsection*{Implicit Regularization and the Bias-Complexity Trade-off}

The capacity control achieved by margin maximization can be understood as a form of
\emph{implicit regularization}. Traditional regularization methods, such as $\ell_2$ weight
decay or $\ell_1$ penalization, add an explicit complexity penalty to the training objective
in order to prevent the learned hypothesis from belonging to an overly expressive class. In
the SVM, this effect arises naturally from the geometry of the problem: minimizing
$\|\mathbf{w}\|^2$ is equivalent to maximizing $\gamma$, which by~\eqref{eq:margin_vc}
minimizes an upper bound on the VC dimension of the effective hypothesis class. No explicit
regularization term is needed.

This perspective also clarifies the role of the hyperparameter $C$ in the soft-margin
formulation introduced earlier,
\begin{myequation*}
\min_{\mathbf{w}, b, \boldsymbol{\xi}} \;\frac{1}{2}\|\mathbf{w}\|^2 + C\sum_{i=1}^n \xi_i
\quad \text{subject to} \quad
t_i(\mathbf{w}^\top \mathbf{x}_i + b) \ge 1 - \xi_i, \;\; \xi_i \ge 0, \quad i=1,\ldots,n.
\end{myequation*}
Large $C$ penalizes margin violations heavily, yielding a narrower margin, a larger effective
VC dimension, and fewer training errors. Small $C$ tolerates more violations in exchange for
a wider margin and lower effective VC dimension. This is precisely the bias-complexity
trade-off that the Fundamental Theorem places at the center of the theory of learning: the
sum of training error and complexity term in~\eqref{eq:svm_gen_bound} must be jointly
minimized, and $C$ is the parameter that mediates this balance. Choosing $C$ by
cross-validation is therefore a principled way to identify the effective capacity
$d_{\mathrm{eff}} = R^2/\gamma^2$ that best controls generalization on the given data.

Finally, this framework gives a clean account of overfitting. Overfitting occurs when the
hypothesis class is too expressive relative to the available data, so that the complexity
term in the VC generalization bound dominates and the generalization gap is large. A small-margin
classifier belongs to a class with large effective VC dimension, leaving it free to exploit
accidental alignments or noise in the training sample. A large-margin classifier, by
contrast, belongs to a class of provably lower VC dimension and is therefore constrained to
produce decision boundaries that are stable under small perturbations of the data --- exactly
the stability that the VC framework identifies as the condition for generalization from finite
samples.

\section{Gradient Descent and  SVM}

%\subsubsection*{Computational issues}
%\begin{itemize}
%  \item Training time of the standard SVM is $O(n^3)$ (solving QP) \begin{itemize}
%  \item Can be prohibitive for large datasets
%  \end{itemize}
%  \item Lots of research has gone into speeding up the SVMs
%\begin{itemize}
%  \item Many approximate QP solvers are used to speed up SVMs
%  \item Gradient descent faster and with possibility of limiting the computation time
%  \end{itemize}
%\end{itemize}
%
%
%The dual formulation of Support Vector Machines provides deep insight into the geometry of the solution and enables the use of kernel methods. At the same time, solving the dual problem requires specialized quadratic programming algorithms, whose computational cost can become prohibitive when the number of training examples is very large. For this reason, it is often advantageous to revisit the primal formulation and ask whether SVMs can be trained using more general-purpose optimization techniques.
%
%An important observation is that, once the constraints are handled appropriately, the soft-margin SVM can be cast as an unconstrained optimization problem. This opens the door to gradient-based methods, including (stochastic) gradient descent, which scale well to large datasets and are widely used in modern machine learning practice.
%
%\medskip
%
%\subsubsection*{From the Constrained Primal to an Unconstrained Objective}
%
%Recall the primal formulation of the soft-margin SVM:
%\begin{myalign}
%\min_{\w,b,\Vxi} \ &\frac{1}{2}\|\w\|^2 + C\sum_{i=1}^{n}\xi_i\\
%\text{subject to} \ &t_i(\w^T\boldsymbol{\phi}(\x_i)+b) \geq 1-\xi_i \quad \text{for } i=1,\ldots,n,\\
%&\xi_i \geq 0 \quad \text{for } i=1,\ldots,n.
%\end{myalign}
%
%At first glance, this is a constrained optimization problem involving both the model parameters $(\w,b)$ and the slack variables $\xi_i$. However, a key simplification arises from the fact that, for fixed $(\w,b)$, the slack variables are completely decoupled from one another and appear only linearly in both the objective and the constraints. As a result, each $\xi_i$ can be optimized independently.
%
%For a given training example $(\x_i,t_i)$, the constraint
%\begin{myequation*}
%t_i(\w^T\boldsymbol{\phi}(\x_i)+b) \geq 1-\xi_i
%\end{myequation*}
%forces $\xi_i$ to be at least $1 - t_i(\w^T\boldsymbol{\phi}(\x_i)+b)$. Since $\xi_i$ is also constrained to be nonnegative and enters the objective with a positive coefficient, the optimal choice is the smallest value compatible with feasibility. This yields the closed-form expression
%\begin{myequation*}
%\xi_i^* = \max_{}\left\{0,\; 1 - t_i(\w^T\boldsymbol{\phi}(\x_i)+b)\right\}.
%\end{myequation*}
%
%Intuitively, if a point already satisfies the margin constraint, no slack is needed. If the constraint is violated, the slack variable measures exactly how far the point lies inside the margin or on the wrong side of the decision boundary.
%
%Substituting these optimal slack values back into the primal objective eliminates the constraints altogether and leads to an unconstrained optimization problem in $(\w,b)$:
%\begin{myequation*}
%J(\w,b) = \frac{1}{2}\|\w\|^2 + C\sum_{i=1}^{n}\max_{}\left\{0,\; 1 - t_i(\w^T\boldsymbol{\phi}(\x_i)+b)\right\}.
%\end{myequation*}
%
%This reformulation makes explicit that training an SVM amounts to minimizing a regularized empirical loss.
\subsection*{Computational Complexity and SVM Optimization}

The dual formulation of Support Vector Machines provides deep insight into the geometry of
the solution and enables the use of kernel methods. At the same time, solving the dual
problem requires specialized quadratic programming algorithms, whose computational cost can
become prohibitive when the number of training examples is very large. The training time of
the standard SVM, based on solving the quadratic programming problem, is of order $O(n^3)$,
where $n$ is the number of training examples. This cubic complexity makes the method
impractical on large datasets, and for this reason a substantial body of research has focused
on developing more efficient methods for SVM training. Two main approaches can be
distinguished. The first consists in the use of approximate solvers for the quadratic
problem, which sacrifice the exact solution in favor of a sufficiently good approximation
obtained in reduced time. The second approach, conceptually different, consists in
reformulating the training problem so as to make it accessible to more general optimization
methods, in particular gradient descent, which scales favorably with the size of the dataset
and offers the possibility of limiting computation time by stopping the optimization before
full convergence.

To pursue this second approach, it is convenient to temporarily set aside the dual
formulation and return to the primal, asking whether the soft-margin SVM can be rewritten
as an unconstrained optimization problem, directly accessible to gradient-based methods.

\subsubsection*{From the Constrained Primal to an Unconstrained Objective}

Recall the primal formulation of the soft-margin SVM:
\begin{equation}
\min_{\w,b,\boldsymbol{\xi}} \quad \frac{1}{2}\|\w\|^2 + C\sum_{i=1}^{n}\xi_i
\end{equation}
subject to the constraints
\begin{equation}
t_i(\w^T\boldsymbol{\phi}(\x_i)+b) \geq 1-\xi_i \quad \text{for } i=1,\ldots,n,
\qquad
\xi_i \geq 0 \quad \text{for } i=1,\ldots,n.
\end{equation}
At first glance, this is a constrained optimization problem involving both the model
parameters $(\w,b)$ and the slack variables $\xi_i$. However, a key simplification arises
from the observation that, for fixed $(\w,b)$, the slack variables are completely decoupled
from one another and appear only linearly in both the objective and the constraints. As a
result, each $\xi_i$ can be optimized independently of the others.

For a given training example $(\x_i, t_i)$, the constraint
\begin{equation}
t_i(\w^T\boldsymbol{\phi}(\x_i)+b) \geq 1-\xi_i
\end{equation}
forces $\xi_i$ to be at least $1 - t_i(\w^T\boldsymbol{\phi}(\x_i)+b)$. Since $\xi_i$ is
also constrained to be non-negative and appears in the objective with positive coefficient
$C$, the optimal choice is the smallest value compatible with feasibility, namely
\begin{equation}
\xi_i^* = \max\left\{0,\; 1 - t_i(\w^T\boldsymbol{\phi}(\x_i)+b)\right\}.
\end{equation}
This expression has a direct geometric interpretation. If the point $\x_i$ already satisfies
the margin constraint, that is, if $t_i(\w^T\boldsymbol{\phi}(\x_i)+b) \geq 1$, then no
slack is needed and $\xi_i^* = 0$. If instead the constraint is violated, the slack variable
measures exactly how far the point lies inside the margin or on the wrong side of the
decision boundary: the larger the violation, the larger the value of $\xi_i^*$.

Substituting these optimal slack values back into the primal objective eliminates the
constraints entirely and yields an unconstrained optimization problem in the parameters
$(\w,b)$ alone:
\begin{equation}
J(\w,b) = \frac{1}{2}\|\w\|^2 + C\sum_{i=1}^{n}\max\left\{0,\; 1 -
t_i(\w^T\boldsymbol{\phi}(\x_i)+b)\right\}.
\end{equation}
This reformulation makes the structure of the SVM training problem explicit: minimizing
$J(\w,b)$ is equivalent to minimizing a regularized empirical loss. The first term,
$\frac{1}{2}\|\w\|^2$, is the regularization term that controls the capacity of the
classifier, as discussed in connection with the VC dimension. The second term is the sum of
the losses on the individual training examples, weighted by the parameter $C$ that balances
the trade-off between margin width and error penalization.

The loss function for a single example,
\begin{equation}
L_{h}(t_i, \w, b) = \max\left\{0,\; 1 - t_i(\w^T\boldsymbol{\phi}(\x_i)+b)\right\},
\end{equation}
is known as the \alert{hinge loss}. It equals zero for correctly classified points that lie
outside the margin, and grows linearly with the degree of margin violation for points that
are misclassified or lie too close to the decision boundary. Compared to the $0$-$1$ loss
used in elementary binary classification, the hinge loss has the crucial advantage of being
convex and, almost everywhere, differentiable, which makes it compatible with
gradient-based optimization methods.

%The unconstrained formulation therefore opens the door to the use of gradient descent for
%SVM training. Since the hinge loss is not differentiable at the point
%$t_i(\w^T\boldsymbol{\phi}(\x_i)+b) = 1$, in practice one uses the \alert{subgradient}, a
%generalization of the gradient applicable to convex but non-differentiable functions. In
%particular, stochastic subgradient descent, which updates the parameters using the gradient
%computed on a single example or a small subset of the dataset at a time, scales favorably
%with $n$ and constitutes one of the most widely used approaches for training SVMs on large
%datasets. 
%%
%
%\subsubsection*{Hinge Loss and Regularized Risk Minimization}
%
%The quantity inside the summation,
%\begin{myequation*}
%\ell_H(\w,b,\x_i,t_i) = \max_{}\left\{0,\; 1 - t_i(\w^T\boldsymbol{\phi}(\x_i)+b)\right\},
%\end{myequation*}
%is known as the \alert{hinge loss}. It assigns zero loss to points that are correctly classified with margin at least one, and a positive loss proportional to the amount of margin violation otherwise.
%
%Geometrically, the hinge loss measures how far a training point lies from the margin boundary. Points that are well separated from the decision boundary play no role in shaping the classifier, whereas points that lie inside the margin or are misclassified exert a corrective influence on the parameters.
%
%With this notation, the SVM objective takes the form
%\begin{myequation*}
%J(\w,b) = \frac{1}{2}\|\w\|^2 + C\sum_{i=1}^{n}\ell_H(\w,b,\x_i,t_i),
%\end{myequation*}
%which is a canonical example of a \alert{regularized loss minimization} problem. The first term controls model complexity by penalizing large weights, while the second term enforces fidelity to the training data. The parameter $C$ governs the trade-off between these two goals: increasing $C$ emphasizes correct classification, whereas decreasing $C$ favors wider margins and stronger regularization.
%
%\medskip
%
%\subsubsection*{Subgradients of the Hinge Loss}

A technical complication arises from the fact that the hinge loss is not differentiable everywhere. In particular, it has a ``kink'' at the margin boundary where
\begin{myequation*}
t_i(\w^T\boldsymbol{\phi}(\x_i)+b) = 1.
\end{myequation*}
Nevertheless, the objective function remains convex, and optimization can proceed using the concept of a \alert{subgradient}, which generalizes the notion of a gradient to nonsmooth convex functions.

%A vector $\mathbf{g}$ is a subgradient of a convex function $f$ at a point $\x$ if
%\begin{myequation*}
%f(\mathbf{y}) \geq f(\x) + \mathbf{g}^T(\mathbf{y}-\x) \quad \text{for all } \mathbf{y}.
%\end{myequation*}
%
%For the hinge loss, a valid subgradient with respect to $\w$ is given by
%\begin{myequation*}
%\frac{\partial \ell_H}{\partial \w} =
%\begin{cases}
%\mathbf{0} & \text{if } t_i(\w^T\boldsymbol{\phi}(\x_i)+b) \geq 1,\\
%-\,t_i\boldsymbol{\phi}(\x_i) & \text{if } t_i(\w^T\boldsymbol{\phi}(\x_i)+b) < 1.
%\end{cases}
%\end{myequation*}
%
%Aggregating over all training points, the subgradient of the full objective with respect to $\w$ takes the form
%\begin{myequation*}
%\frac{\partial J}{\partial \w} = \w - C\sum_{i \in V} t_i\boldsymbol{\phi}(\x_i),
%\end{myequation*}
%where $V$ denotes the set of indices corresponding to margin-violating examples, that is,
%\begin{myequation*}
%V = \left\{ i \mid t_i(\w^T\boldsymbol{\phi}(\x_i)+b) < 1 \right\}.
%\end{myequation*}
%
%This expression makes clear that only points violating the margin constraint contribute to the data-dependent part of the update.
%
%\medskip
%
%\subsubsection*{Gradient Descent Updates}
%
%Once the problem has been cast in unconstrained form, standard gradient-based methods can be applied. In batch gradient descent, all parameters are updated simultaneously using the full subgradient at the current iterate:
%\begin{myequation*}
%\w^{(t+1)} = \w^{(t)} - \eta \frac{\partial J}{\partial \w}\bigg|_{\w=\w^{(t)}},
%\end{myequation*}
%where $\eta>0$ is the learning rate.
%
%Substituting the expression for the subgradient yields
%\begin{myequation*}
%\w^{(t+1)} = \w^{(t)} - \eta\left(\w^{(t)} - C\sum_{i \in V^{(t)}} t_i\boldsymbol{\phi}(\x_i)\right),
%\end{myequation*}
%or, equivalently,
%\begin{myequation*}
%\w^{(t+1)} = (1-\eta)\w^{(t)} + \eta C\sum_{i \in V^{(t)}} t_i\boldsymbol{\phi}(\x_i).
%\end{myequation*}
%
%This update admits a clear interpretation. The multiplicative factor $(1-\eta)$ shrinks the current weight vector, implementing the effect of regularization and discouraging overly complex models. The additive term pushes the solution in a direction that corrects margin violations, reinforcing the influence of misclassified or weakly classified points. Training examples that are correctly classified with sufficient margin do not affect the update at all.
%
%The bias parameter $b$ can be updated in an analogous way using subgradients, or treated separately depending on the implementation.
\subsubsection*{Subgradient of the Hinge Loss}

To apply gradient-based methods to the minimization of the unconstrained objective
$J(\w,b)$, it is necessary to compute the derivatives of the hinge loss with respect to
the model parameters. Since the hinge loss is not differentiable at the point
$t_i(\w^T\boldsymbol{\phi}(\x_i)+b) = 1$, one must resort to the concept of
\alert{subgradient}, which generalizes the notion of gradient to convex but
non-differentiable functions.

Formally, a vector $\mathbf{g}$ is a subgradient of a convex function $f$ at a point $\x$
if it satisfies the inequality
\begin{myequation*}
f(\mathbf{y}) \geq f(\x) + \mathbf{g}^T(\mathbf{y}-\x) \quad \text{for all } \mathbf{y}.
\end{myequation*}
This condition generalizes the inequality that characterizes differentiable convex functions:
a subgradient is any vector that defines a supporting hyperplane to the graph of $f$ at the
point $\x$. When $f$ is differentiable, the unique subgradient coincides with the ordinary
gradient; at points of non-differentiability, multiple subgradients may exist, and any valid
choice can be used in the optimization updates.

For the hinge loss $\ell_H(\w) = \max\{0, 1 - t_i(\w^T\boldsymbol{\phi}(\x_i)+b)\}$, a
valid subgradient with respect to $\w$ is given by
\begin{myequation*}
\frac{\partial \ell_H}{\partial \w} =
\begin{cases}
\mathbf{0} & \text{if } t_i(\w^T\boldsymbol{\phi}(\x_i)+b) \geq 1,\\[4pt]
-\,t_i\boldsymbol{\phi}(\x_i) & \text{if } t_i(\w^T\boldsymbol{\phi}(\x_i)+b) < 1.
\end{cases}
\end{myequation*}
The interpretation is immediate: if the point $\x_i$ is correctly classified with sufficient
margin, the hinge loss is zero and its subgradient is the zero vector, meaning the point
does not contribute to the update. If instead the margin constraint is violated --- the
point is misclassified or lies too close to the decision boundary --- the subgradient is
$-t_i\boldsymbol{\phi}(\x_i)$, a term that depends on the label and on the representation
of the point in the feature space.

Aggregating the contributions of all training examples, the subgradient of the full
objective $J(\w,b)$ with respect to $\w$ is
\begin{myequation*}
\frac{\partial J}{\partial \w} = \w - C\sum_{i \in V} t_i\boldsymbol{\phi}(\x_i),
\end{myequation*}
where $V$ is the set of indices of examples that violate the margin constraint,
\begin{myequation*}
V = \left\{ i \;\middle|\; t_i(\w^T\boldsymbol{\phi}(\x_i)+b) < 1 \right\}.
\end{myequation*}
This expression highlights an important structural property: the data-dependent part of the
subgradient receives contributions exclusively from the examples in $V$, that is, from those
that lie inside the margin or on the wrong side of the decision boundary. Examples that are
correctly classified with sufficient margin do not influence the update at all. This property
is analogous, in this context, to what was already observed in the dual problem: there too,
only the support vectors --- the points lying exactly on the margin --- contributed to the
solution. Here, during iterative optimization, it is the points that violate the current
margin that drive the updates.

\subsubsection*{Gradient Descent Updates}

Once the problem has been cast in unconstrained form, standard gradient-based optimization
methods can be applied. In \alert{batch gradient descent}, all parameters are updated
simultaneously using the full subgradient computed over the entire dataset at the current
iterate:
\begin{myequation*}
\w^{(t+1)} = \w^{(t)} - \eta \frac{\partial J}{\partial \w}\bigg|_{\w=\w^{(t)}},
\end{myequation*}
where $\eta > 0$ is the learning rate, a hyperparameter that controls the step size of each
update. Substituting the expression for the subgradient yields
\begin{myequation*}
\w^{(t+1)} = \w^{(t)} - \eta\left(\w^{(t)} - C\sum_{i \in V^{(t)}}
t_i\boldsymbol{\phi}(\x_i)\right),
\end{myequation*}
which can be rewritten in the equivalent form
\begin{myequation*}
\w^{(t+1)} = (1-\eta)\w^{(t)} + \eta C\sum_{i \in V^{(t)}} t_i\boldsymbol{\phi}(\x_i).
\end{myequation*}
This form of the update admits a particularly clear interpretation. The multiplicative
factor $(1-\eta)$, assuming $\eta < 1$, contracts the current weight vector toward the
origin at each iteration. This contraction effect directly implements the role of the
regularization term $\frac{1}{2}\|\w\|^2$ in the objective: without the correction provided
by the second term, the weights would decay geometrically toward zero, discouraging
excessively complex models. The additive term
$\eta C\sum_{i \in V^{(t)}} t_i\boldsymbol{\phi}(\x_i)$ corrects this contraction by
pushing the solution in the direction that reduces margin violations: for each example
$i \in V^{(t)}$, the contribution $t_i\boldsymbol{\phi}(\x_i)$ orients the update so as
to move the classifier away from the decision boundary in the correct direction for that
point. Examples that are correctly classified with sufficient margin do not appear in
$V^{(t)}$ and have no influence on the update whatsoever.

The parameter $C$ mediates the balance between these two opposing effects: a large value
of $C$ amplifies the corrective contribution of the violating examples, reducing the
emphasis on regularization and favoring more accurate classification on the training set;
a small value of $C$ instead favors contraction of the weights and a wider margin, even at
the cost of tolerating more violations. This is the same tension between training error and
model complexity already identified in the context of the VC dimension and the Fundamental
Theorem of Statistical Learning.

The bias parameter $b$ can be updated in an analogous way using its subgradient, or treated
separately depending on the implementation. In many practical implementations, $b$ is not
subject to regularization, which is reflected in an update that lacks the contraction term.

In practice, for large datasets, the stochastic variant of subgradient descent is preferred
over the batch version: at each iteration, the subgradient is computed on a single example
or a small subset of the dataset, drastically reducing the cost per iteration at the expense
of greater variability in the updates. This variability is often advantageous, as it allows
the parameter space to be explored more freely and tends to converge more quickly toward a
solution of good quality.

\medskip

All other variants of gradient descent, including stochastic and mini-batch methods, can be derived in the same spirit by replacing the full sum over $V$ with estimates based on individual examples or small subsets of the data. These methods dramatically reduce the computational cost per iteration and are the standard choice for large-scale learning problems.


\section*{Kernel Methods: Enabling Nonlinear Learning}

Up to this point, we have emphasized that many learning algorithms---including linear regression and linear SVMs---are attractive because they are computationally efficient, convex, and relatively easy to analyze. Their main limitation is also clear: a linear decision boundary (or a linear predictor) can only capture linear patterns in the input variables. When the underlying structure of the data is genuinely nonlinear, linear models can perform poorly even when regularized and carefully tuned.

Kernel methods provide a principled way to overcome this limitation without abandoning the advantages of linear methods. The central idea is to keep the learning algorithm \emph{linear} in a suitable representation of the data, while allowing the \emph{decision boundary in the original input space} to become nonlinear. This is achieved by working in a transformed feature space, but doing so implicitly through inner products, avoiding explicit construction of high-dimensional feature vectors. This mechanism is often referred to as the \alert{kernel trick}.

\medskip

\subsubsection*{Motivation: Why Nonlinear Models Matter}

Nonlinearity arises naturally in both classification and regression. In classification, the classes may not be separable by any hyperplane in the original space; classic examples include concentric circles or spiral-shaped clusters, where any linear boundary necessarily makes many errors. In regression, the relationship between inputs and outputs may follow a curved trend, include oscillations, or depend on interactions between variables that cannot be captured by a linear function.

A naive approach to nonlinearity is to manually enrich the representation: add polynomial terms, cross-products, trigonometric features, or other nonlinear transformations. While this idea is conceptually simple, it quickly becomes impractical. First, generating a large set of candidate nonlinear features can be computationally expensive. Second, deciding which features to include often requires substantial domain expertise. Third, the dimension of the expanded space can explode, creating serious memory and runtime bottlenecks and increasing the risk of overfitting.

Kernel methods address exactly these issues: they allow us to behave \emph{as if} we had constructed a rich nonlinear feature representation, but without explicitly forming that representation.

\medskip

\subsubsection*{Feature Maps and Kernel Functions: Formal Setup}

The formal starting point is a basis function, or \alert{feature map} 
\begin{myequation*}
\boldsymbol{\phi}: \mathcal{X} \to \mathcal{F},
\end{myequation*}
where $\mathcal{X}$ is the original input space (typically $\mathbb{R}^d$) and $\mathcal{F}$ is a (possibly very high-dimensional) \alert{feature space}. The space $\mathcal{F}$ is assumed to be a \alert{Hilbert space}, meaning it is a vector space equipped with an inner product (and which satisfies  a completeness condition that we will ignore here for simplicity of exposition). 

Once $\boldsymbol{\phi}$ is fixed, we apply a linear model in feature space. For SVMs, this means considering decision functions of the form
\begin{myequation*}
f(\x) = \w^T\boldsymbol{\phi}(\x) + b.
\end{myequation*}
The decision boundary is linear in $\mathcal{F}$, but it can correspond to a highly nonlinear surface when viewed back in the original space $\mathcal{X}$.

The difficulty is that explicitly computing $\boldsymbol{\phi}(\x)$ may be expensive or even impossible if $\mathcal{F}$ is infinite-dimensional. Kernel methods avoid this by noticing that many learning algorithms---and SVMs in particular---can be expressed entirely in terms of inner products between feature vectors. This motivates the definition of a \alert{kernel function}
\begin{myequation*}
\kappa(\x_i, \x_j) = \boldsymbol{\phi}(\x_i)^T\boldsymbol{\phi}(\x_j).
\end{myequation*}
A kernel computes the inner product in feature space while taking as input only the original vectors $\x_i$ and $\x_j$. The feature vectors themselves never need to be constructed. This identity is the essence of the kernel trick: wherever an algorithm requires $\boldsymbol{\phi}(\x_i)^T\boldsymbol{\phi}(\x_j)$, we replace it with $\kappa(\x_i,\x_j)$ and remain entirely in the original input space.

\medskip

\subsubsection*{Valid Kernels: Mercer's Condition}

Not every function $\kappa(\x,\x')$ can be used as a kernel. For the kernel trick to be meaningful, the function must behave like an inner product in some feature space. The standard characterization is expressed through positive semidefiniteness.

\medskip

\noindent
\textbf{Definition (positive semidefinite kernel).}
A function $\kappa: \mathcal{X}\times\mathcal{X}\to\mathbb{R}$ is a valid kernel if and only if, for every finite set of points $\{\x_1,\ldots,\x_n\}\subseteq\mathcal{X}$, the associated \alert{Gram matrix} $\mathbf{K}$ defined by $K_{ij}=\kappa(\x_i,\x_j)$ is positive semidefinite, namely
\begin{myequation*}
\mathbf{v}^T\mathbf{K}\mathbf{v} \ge 0 \quad \text{for all } \mathbf{v}\in\mathbb{R}^n.
\end{myequation*}
Equivalently, all eigenvalues of $\mathbf{K}$ are nonnegative.

\medskip

This condition guarantees that there exists at least one Hilbert space $\mathcal{F}$ and a feature map $\boldsymbol{\phi}$ such that $\kappa(\x,\x')$ equals the inner product of $\boldsymbol{\phi}(\x)$ and $\boldsymbol{\phi}(\x')$. In practical terms, positive semidefiniteness is what makes kernels safe to use: it ensures that the implicit geometry induced by $\kappa$ is consistent and that optimization problems such as the SVM dual are well-posed.

\medskip

\subsubsection*{Constructing Kernels: Techniques}

In practice, one rarely verifies the positive semidefinite condition from scratch for an arbitrary function. Instead, kernels are typically constructed from simpler kernels using operations that preserve validity. If $\kappa_1$ and $\kappa_2$ are valid kernels, then the following constructions produce valid kernels as well:
\begin{enumerate}
\item \textbf{Sum:} $\kappa(\x_1,\x_2) = \kappa_1(\x_1,\x_2) + \kappa_2(\x_1,\x_2)$
\item \textbf{Product:} $\kappa(\x_1,\x_2) = \kappa_1(\x_1,\x_2)\cdot \kappa_2(\x_1,\x_2)$
\item \textbf{Scalar multiple:} $\kappa(\x_1,\x_2) = c\,\kappa_1(\x_1,\x_2)$ for any $c>0$
\item \textbf{Exponential:} $\kappa(\x_1,\x_2) = e^{\kappa_1(\x_1,\x_2)}$
\item \textbf{Polynomial:} $\kappa(\x_1,\x_2) = p(\kappa_1(\x_1,\x_2))$ for any polynomial $p$ with nonnegative coefficients
\item \textbf{Weighted inner product:} $\kappa(\x_1,\x_2) = \x_1^T\mathbf{A}\x_2$ for any positive definite matrix $\mathbf{A}$
\item \textbf{Scaling:} $\kappa(\x_1,\x_2) = f(\x_1)\kappa_1(\x_1,\x_2)g(\x_2)$ for any functions $f,g:\mathcal{X}\to\mathbb{R}$
\item \textbf{Feature map composition:} $\kappa(\x_1,\x_2) = \kappa_3(\boldsymbol{\phi}(\x_1),\boldsymbol{\phi}(\x_2))$ where $\boldsymbol{\phi}$ is any feature map and $\kappa_3$ is any kernel on the transformed space
\end{enumerate}

These closure properties explain why kernel design is flexible: complex similarity measures can be engineered by combining simpler building blocks, all while guaranteeing that the resulting kernel corresponds to some valid feature space geometry.

\medskip

\subsubsection*{Polynomial Kernels: Example Construction}

A canonical example is the polynomial kernel
\begin{myequation*}
\kappa(\x_1,\x_2) = (\x_1^T\x_2 + c)^d.
\end{myequation*}

This kernel can be justified constructively. The basic dot product $\x_1^T\x_2$ is itself a kernel, corresponding to the identity feature map. A constant kernel $c$ is also valid, corresponding to adding a constant feature. Their sum $\x_1^T\x_2+c$ is therefore a kernel. Finally, raising a kernel to an integer power $d$ (equivalently applying the polynomial $p(z)=z^d$ with nonnegative coefficients) produces another valid kernel.

To make the feature expansion explicit, consider $d=2$ in $\mathbb{R}^2$. We obtain
\begin{myalign}
\kappa(\x_1,\x_2) &= (x_{1,1}x_{2,1} + x_{1,2}x_{2,2} + c)^2\\
&= \sum_{i,j} x_{1,i}x_{1,j}x_{2,i}x_{2,j} + 2c\sum_i x_{1,i}x_{2,i} + c^2.
\end{myalign}
This corresponds to the feature map
\begin{myequation*}
\boldsymbol{\phi}(\x) =
(x_1^2,\ x_2^2,\ \sqrt{2}x_1x_2,\ \sqrt{2c}x_1,\ \sqrt{2c}x_2,\ c)^T,
\end{myequation*}
which maps from $\mathbb{R}^2$ to $\mathbb{R}^6$.

The computational point is crucial. Even though the feature dimension grows rapidly with the polynomial degree, evaluating $\kappa(\x_1,\x_2)$ still requires only $O(d)$ work: a single dot product followed by a scalar power. In contrast, explicitly constructing the polynomial feature vector of degree $t$ would require a number of components on the order of $O(d^t)$. The kernel trick therefore provides a dramatic savings in computation and memory.

\medskip

\subsubsection*{Gaussian Kernel: Infinite-Dimensional Features}

Among the most widely used kernels is the \alert{Gaussian} (or RBF) kernel
\begin{myequation*}
\kappa(\x_1,\x_2) = \exp\left(-\frac{\|\x_1 - \x_2\|^2}{2\sigma^2}\right),
\end{myequation*}
where $\sigma>0$ is a bandwidth parameter.

To see why this function is a valid kernel, one can start from the identity
\begin{myequation*}
\|\x_1 - \x_2\|^2 = \|\x_1\|^2 + \|\x_2\|^2 - 2\x_1^T\x_2
\end{myequation*}
and rewrite the exponential as
\begin{myequation*}
\kappa(\x_1,\x_2)
= e^{-\|\x_1\|^2/(2\sigma^2)} \cdot e^{\x_1^T\x_2/\sigma^2} \cdot e^{-\|\x_2\|^2/(2\sigma^2)}.
\end{myequation*}

The middle factor is the exponential of a scaled linear kernel and defines a valid kernel, while the outer factors depend only on $\x_1$ and $\x_2$ separately and can be seen as scaling terms; scaling a valid kernel in this way preserves validity. This decomposition provides an explicit route to verifying that the Gaussian kernel satisfies the kernel conditions.

\medskip

The feature-space interpretation is equally important: the Gaussian kernel corresponds to an infinite-dimensional feature map. Although $\boldsymbol{\phi}(\x)$ has infinitely many components, the kernel trick allows inner products to be computed exactly with a simple closed-form expression. This is the source of the Gaussian kernel's remarkable flexibility: it can represent extremely rich families of decision boundaries while keeping computations manageable through kernel evaluations.

The bandwidth $\sigma$ controls the locality of the similarity measure. When $\sigma$ is small, $\kappa(\x_1,\x_2)$ is close to $1$ only for very nearby points, which allows highly flexible, finely detailed decision boundaries but increases the risk of overfitting. When $\sigma$ is large, the kernel becomes smoother and more global, which can improve robustness but may lead to underfitting. As with $C$ in soft-margin SVMs, $\sigma$ is typically selected by cross-validation.

\subsubsection*{Common Kernel Functions and When to Use Them}

Different kernels encode different inductive biases, meaning different assumptions about what kinds of patterns are likely to appear in the data. The most common choices include:

\medskip

\noindent
\textbf{Linear kernel.} The simplest kernel is
\begin{myequation*}
\kappa(\x_1,\x_2) = \x_1^T\x_2.
\end{myequation*}
It is appropriate when the data are already (approximately) linearly separable and is the fastest and least prone to overfitting. It also preserves interpretability in the original feature space.

\medskip

\noindent
\textbf{Polynomial kernel.} A standard nonlinear alternative is
\begin{myequation*}
\kappa(\x_1,\x_2) = (\x_1^T\x_2 + c)^d.
\end{myequation*}
It introduces controlled nonlinearity: the degree $d$ governs complexity, while $c$ adjusts the relative weight of lower- and higher-order terms. Polynomial kernels are often effective when interactions between features matter but extremely localized boundaries are not needed.

\medskip

\noindent
\textbf{Gaussian (RBF) kernel.} The Gaussian kernel can also be written as
\begin{myequation*}
\kappa(\x_1,\x_2) = \exp(-\gamma\|\x_1 - \x_2\|^2),
\end{myequation*}
where $\gamma = 1/(2\sigma^2)$. It is a strong default choice when one expects complex nonlinear structure and has enough data to tune hyperparameters and avoid overfitting. Its performance can degrade in very high dimensions due to distance concentration effects.

\medskip

\noindent
\textbf{Sigmoid kernel.} Another option is
\begin{myequation*}
\kappa(\x_1,\x_2) = \tanh(c_1 \x_1^T\x_2 + c_2).
\end{myequation*}
It is historically motivated by neural networks, but it does not satisfy the kernel validity condition for all parameter choices, so it should be used with care.

\medskip

\subsection*{Kernels in SVM: The Full Picture}

We can now recall and summarize how kernels enter the SVM training and prediction pipeline. In the soft-margin case, the dual optimization problem takes the kernelized form
\begin{myalign}
&\max_{\Vlambda} \left(\sum_{i=1}^{n}\lambda_i - \frac{1}{2}\sum_{i=1}^{n}\sum_{j=1}^{n}\lambda_i\lambda_j t_i t_j \kappa(\x_i, \x_j)\right)\\
&\text{subject to} \quad 0 \leq \lambda_i \leq C \quad \text{for } i=1,\ldots,n\\
&\quad\quad\quad\quad \sum_{i=1}^{n}\lambda_i t_i = 0.
\end{myalign}

Once the optimal multipliers $\Vlambda^*$ are found, prediction for a new point $\x$ uses the decision function
\begin{myequation*}
y(\x) = \sum_{i \in \mathcal{S}}\lambda_i^* t_i \kappa(\x_i, \x) + b^*,
\end{myequation*}
where $\mathcal{S}$ is the set of support vectors.

This representation makes clear why kernelized SVMs are powerful. Nonlinearity is entirely contained in the kernel $\kappa$, while the learning algorithm remains a convex optimization problem with strong theoretical guarantees. At the same time, sparsity is preserved: only support vectors contribute to predictions, which keeps evaluation efficient. Finally, kernel choice provides a flexible way to inject prior knowledge about similarity and structure into the learning process, enabling SVMs to adapt to a wide variety of datasets without changing the underlying optimization framework.



%\end{document}